\section{AIME II 2023}
        \begin{enumerate}
        	\item (\textit{AIME II 2023}) The numbers of apples growing on each of six apple trees form an arithmetic sequence where the greatest number of apples growing on any of the six trees is double the least number of apples growing on any of the six trees. The total number of apples growing on all six trees is $990.$ Find the greatest number of apples growing on any of the six trees.


 
 \textbf{Answer: }220



\textbf{Solution 1}

In the arithmetic sequence, let $a$ be the first term and $d$ be the common difference, where $d>0.$ The sum of the first six terms is 
 \[a+(a+d)+(a+2d)+(a+3d)+(a+4d)+(a+5d) = 6a+15d.\]We are given that
 \begin{align*} 6a+15d &= 990, \\ 2a &= a+5d. \end{align*}
The second equation implies that $a=5d.$ Substituting this into the first equation, we get 
 \begin{align*} 6(5d)+15d &=990, \\ 45d &= 990 \\ d &= 22. \end{align*}
It follows that $a=110.$ Therefore, the greatest number of apples growing on any of the six trees is $a+5d=\boxed{220}.$




\textbf{Solution 2}

Let the terms in the sequence be defined as 
 \[a_1, a_2, ..., a_6.\] 

 Since this is an arithmetic sequence, we have $a_1+a_6=a_2+a_5=a_3+a_4.$ So, 
 \[\sum_{i=1}^6 a_i=3(a_1+a_6)=990.\] Hence, $(a_1+a_6)=330.$ And, since we are given that $a_6=2a_1,$ we get $3a_1=330\implies a_1=110$ and $a_6=\boxed{220}.$



	\item (\textit{AIME II 2023}) Recall that a palindrome is a number that reads the same forward and backward. Find the greatest integer less than $1000$ that is a palindrome both when written in base ten and when written in base eight, such as $292 = 444_{\text{eight}}.$

 \textbf{Answer: }585



\textbf{Solution}

Assuming that such palindrome is greater than $777_8 = 511,$ we conclude that the palindrome has four digits when written in base $8.$ Let such palindrome be 
 \[(\underline{ABBA})_8 = 512A + 64B + 8B + A = 513A + 72B.\]

 It is clear that $A=1,$ so we repeatedly add $72$ to $513$ until we get palindromes less than $1000:$
 \begin{align*} 513+72\cdot0 &= 513, \\ 513+72\cdot1 &= \boxed{585}, \\ 513+72\cdot2 &= 657, \\ 513+72\cdot3 &= 729, \\ 513+72\cdot4 &= 801, \\ 513+72\cdot5 &= 873, \\ 513+72\cdot6 &= 945, \\ 513+72\cdot7 &= 1017. \\ \end{align*}





\textbf{Solution 2 (Bash)}

$1000$ in base $8$ is $1750$, That means any palindrome less has to ends and begin with a 1. You test options until you get to $1111_{8}=585$ which is the answer.

 


	\item (\textit{AIME II 2023}) Let $\triangle ABC$ be an isosceles triangle with $\angle A = 90^\circ.$ There exists a point $P$ inside $\triangle ABC$ such that $\angle PAB = \angle PBC = \angle PCA$ and $AP = 10.$ Find the area of $\triangle ABC.$

 \textbf{Answer: }250



\textbf{Solution 1}

This solution refers to the \textbf{Diagram} section.

 Let $\angle PAB = \angle PBC = \angle PCA = \theta,$ from which $\angle PAC = 90^\circ-\theta,$ and $\angle APC = 90^\circ.$

 Moreover, we have $\angle PBA = \angle PCB = 45^\circ-\theta,$ as shown below:
                    \begin{figure}[H]
                        \centering
                        \includegraphics[width=0.5\textwidth]{Solution/AIME_II_2023_3_Solution_Figure_1.png}
                    \end{figure}
                    Note that $\triangle PAB \sim \triangle PBC$ by AA Similarity. The ratio of similitude is $\frac{PA}{PB} = \frac{PB}{PC} = \frac{AB}{BC},$ so $\frac{10}{PB} = \frac{1}{\sqrt2}$ and thus $PB=10\sqrt2.$ Similarly, we can figure out that $PC=20$.

 Finally, $AC=\sqrt{10^2+20^2}=10\sqrt{5}$, so the area of $\triangle ABC$ is 
 \[\frac12\cdot AB\cdot AC = \frac12\cdot (10\sqrt{5})^2 = \boxed{250}.\]




\textbf{Solution 2}

Since the triangle is a right isosceles triangle, $\angle B = \angle C = 45^\circ$.

 Let the common angle be $\theta$. Note that $\angle PAC = 90^\circ-\theta$, thus $\angle APC = 90^\circ$. From there, we know that $AC = \frac{10}{\sin\theta}$.

 Note that $\angle ABP = 45^\circ-\theta$, so from law of sines we have
 \[\frac{10}{\sin\theta \cdot \frac{\sqrt{2}}{2}}=\frac{10}{\sin(45^\circ-\theta)}.\]Dividing by $10$ and multiplying across yields
 \[\sqrt{2}\sin(45^\circ-\theta)=\sin\theta.\]From here use the sine subtraction formula, and solve for $\sin\theta$:
 \begin{align*} \cos\theta-\sin\theta&=\sin\theta \\ 2\sin\theta&=\cos\theta \\ 4\sin^2\theta&=\cos^2\theta \\ 4\sin^2\theta&=1-\sin^2\theta \\ 5\sin^2\theta&=1 \\ \sin\theta&=\frac{1}{\sqrt{5}}. \end{align*}
Substitute this to find that $AC=10\sqrt{5}$, thus the area is $\frac{(10\sqrt{5})^2}{2}=\boxed{250}$.

 



\textbf{Solution 3}

Since the triangle is a right isosceles triangle, $\angle B = \angle C = 45^\circ$.

 Do some angle chasing yielding:

 
            \begin{itemize}
        \item $\angle APB = \angle BPC = 135^\circ$
            \end{itemize}
            
            \begin{itemize}
        \item $\angle APC=90^\circ$
            \end{itemize}
            We have $AC=\frac{10}{\sin\theta}$ since $\triangle APC$ is a right triangle. Since $\triangle ABC$ is a $45^\circ$-$45^\circ$-$90^\circ$ triangle, $AB=\frac{10}{\sin\theta}$, and $BC=\frac{10\sqrt{2}}{\sin\theta}$. 

 Note that $\triangle APB \sim \triangle BPC$ by a factor of $\sqrt{2}$. Thus, $BP = 10\sqrt{2}$, and $PC = 20$.

 From Pythagorean theorem, $AC=10\sqrt{5}$ so the area of $\triangle ABC$ is $\frac{(10\sqrt{5})^2}{2}=\boxed{250}$.

 



\textbf{Solution 4}

Since the triangle is a right isosceles triangle, $\angle B = \angle C = 45^\circ$.

 Notice that in triangle $PBC$, $\angle PBC + 45-\angle PCA = 45^\circ$, so $\angle BPC = 135^\circ$. Similar logic shows $\angle APC = 135^\circ$.

 Now, we see that $\triangle APB \sim \triangle BPC$ with ratio $1:\sqrt{2}$ (as $\triangle ABC$ is a $45^\circ$-$45^\circ$-$90^\circ$ triangle). Hence, $\overline{PB}=10\sqrt{2}$. We use the Law of Cosines to find $AB$.
 \begin{align*} AB^2&=BP^2+AP^2-2ab\cos(APB) \\ &=100+200-2(10)(10\sqrt{2}\cos(135^\circ)) \\ &=300+200\cdot\sqrt{2}\cdot\frac{1}{\sqrt{2}} \\ &=500. \end{align*}
Since $\triangle ABC$ is a right triangle, the area is $\frac{AB^2}{2}=\frac{500}{2}=\boxed{250}$.

 



\textbf{Solution 5}

Denote the area of $X$ by $[X].$ As in previous solutions, we see that $\angle APC = 90 ^\circ, \triangle BPC \sim \triangle APB$ with ratio $k = \sqrt{2}\implies$ 
 \[\frac {PC}{PB} = \frac {PB}{PA} = k \implies PC = k^2 \cdot AP = 20 \implies [APC] = \frac {AP \cdot PC}{2} = 100.\]
 \[[BPC] = k^2 [APB] = 2 [APB].\]
 \[AB = BC, \angle PCA = \angle PAB \implies \frac {[APC]}{[APB]} = \frac {PC}{PA} = 2 \implies\] 
 \[[ABC] = [APB] + [APC] + [BPC] = [APC] \cdot (\frac {1}{2} + 1 + 2 \cdot \frac {1}{2}) = \frac {5}{2} \cdot  [APC] = \boxed{250}.\]\textbf{vladimir.shelomovskii@gmail.com, vvsss}

 



\textbf{Solution 6}

Denote $\angle PCA = \theta$. Then, by trig Ceva's: 
 \begin{align*} \frac{\sin^3(\theta)}{\sin(90-\theta) \cdot \left(\sin(45-\theta)\right)^2} &= 1 \\ \sin^3(\theta) &= \cos(\theta) \cdot \left(\sin(45) \cos(\theta) - \cos(45) \sin(\theta)\right)^2 \\ 2\sin^3(\theta) &= \cos(\theta) \cdot \left(\cos(\theta) - \sin(\theta)\right)^2 \\ 2\sin^2(\theta) &= \cot(\theta) \cdot \left(1 - 2\sin(\theta)\cos(\theta)\right) \\ 2\sin^2(\theta) &= \cot(\theta) - 2\cos^2(\theta) \\ \cot(\theta) &= 2 \\ \sin(\theta) &= \frac{\sqrt{5}}{5}. \end{align*}
Note that $\angle APC$ is a right angle. Therefore: 
 \begin{align*} \sin(\theta) &= \frac{AP}{AC} \\ AC &= \frac{10}{\frac{\sqrt{5}}{5}} \\ &= 10\sqrt{5} \\ |ABC| &= \frac{AC^2}{2} \\ &= \boxed{250}. \end{align*}
~ConcaveTriangle

 



\textbf{Solution 7}

Notice that point $P$ is one of the two Brocard Points of $\triangle ABC$. (The angle equalities given in the problem are equivalent to the definition of a Brocard point.) By the Brocard point formula,
 \begin{align*} \cot(\phi) = \cot(A)+\cot(B)+\cot(C)  \end{align*}
, where $\phi$ is equal to $\angle PAB$.(This is also called the Brocard angle of triangle ABC). Because the triangle is an isosceles right triangle, the cotangents are easy to compute:
 \begin{align*}\cot(\phi) = 0 + 1 + 1 \\ \cot(\phi) = 2\end{align*}
By definition, $\cot(\phi) = \frac{\cos(\phi)}{\sin(\phi)}$. By the Pythagorean identity, $\cos(\phi)=\frac{2\sqrt{5}}{5}$ and $\sin(\phi) = \cos(\phi)=\frac{\sqrt{5}}{5}$. Consider triangle $APB$. By the problem condition, $\angle PBA = 45-\phi$, so $\angle BPA = 135^{\circ}$
 \begin{align*}\sin{45-\theta} = \sin{45}\cos{\phi}-\cos{45}\sin{\phi} = \frac{\sqrt{10}}{10}\end{align*}
Now, we can use the Law of Sines.
 \begin{align*} \frac{AP}{\sin{45-\theta}}&=\frac{AB}{\sin{135}} \\ 10 \sqrt{10} &= \sqrt{2} AB \\ AB &= 10 \sqrt{5} \end{align*}
Therefore, the answer is 
 \[[ABC] = \frac{1}{2} (AB)^2 = \boxed{250}.\]~ewei12

 



\textbf{Solution 8}

Notice that $\angle APC = 90^{\circ}$, $\angle BPA = 135^{\circ}$, and $\angle CPB = 135^{\circ}$ (from Solution 4). Now let $a = 0$ and $p = 10i$. Then by the angle restrictions $c = m + 10i$ and $b = -n + (n+10)i$ for some $m, n$. Since $\angle BAC = 90^{\circ}$, $ci = b$, or $(m + 10i)i = -n + (n+10)i$. Therefore $n = 10$, $AB = \sqrt{10^2 + 20^2} = \sqrt{500}$, and $[ABC] = \frac{1}{2} AB^2 = \boxed{250}$.

 


	\item (\textit{AIME II 2023}) Let $x,y,$ and $z$ be real numbers satisfying the system of equations

 \begin{align*} xy + 4z &= 60 \\ yz + 4x &= 60 \\ zx + 4y &= 60. \end{align*}

Let $S$ be the set of possible values of $x.$ Find the sum of the squares of the elements of $S.$

 \textbf{Answer: }273



\textbf{Solution 1}

We first subtract the second equation from the first, noting that they both equal $60$. 
 \begin{align*} xy+4z-yz-4x&=0 \\ 4(z-x)-y(z-x)&=0 \\ (z-x)(4-y)&=0 \end{align*}


 Case 1: Let $y=4$.

 The first and third equations simplify to:
 \begin{align*} x+z&=15 \\ xz&=44 \end{align*}
from which it is apparent that $x=4$ and $x=11$ are solutions.

 Case 2: Let $x=z$.

 The first and third equations simplify to:
 \begin{align*} xy+4x&=60 \\ x^2+4y&=60 \end{align*}


 We subtract the following equations, yielding:
 \begin{align*} x^2+4y-xy-4x&=0 \\ x(x-4)-y(x-4)&=0 \\ (x-4)(x-y)&=0 \end{align*}


 We thus have $x=4$ and $x=y$, substituting in $x=y=z$ and solving yields $x=6$ and $x=-10$.

 Then, we just add the squares of the solutions (make sure not to double count the $4$), and get 
 \[4^2+11^2+6^2+(-10)^2=16+121+36+100=\boxed{273}.\]~SAHANWIJETUNGA

 



\textbf{Solution 2}

We index these equations as (1), (2), and (3), respectively.Taking $(1)-(2)$, we get
 \[ \left( x - z \right) \left( y - 4 \right) = 0 . \]

 Denote $x' = x - 4$, $y' = y - 4$, $z' = z - 4$.Thus, the above equation can be equivalently written as
 \[ \left( x' - z' \right) y' = 0 . \hspace{1cm} (1') \]

 Similarly, by taking $(2)-(3)$, we get
 \[ \left( y' - x' \right) z' = 0 . \hspace{1cm} (2') \]

 By taking $(3) - (1)$, we get
 \[ \left( z' - y' \right) x' = 0 . \hspace{1cm} (3') \]

 From $(3')$, we have the following two cases.

 Case 1: $x' = 0$.

 Plugging this into $(1')$ and $(2')$, we get $y'z' = 0$.Thus, $y' = 0$ or $z' = 0$.

 Because we only need to compute all possible values of $x$, without loss of generality, we only need to analyze one case that $y' = 0$.

 Plugging $x' = 0$ and $y' = 0$ into (1), we get a feasible solution $x = 4$, $y = 4$, $z = 11$.

 Case 2: $x' \neq 0$ and $z' - y' = 0$.

 Plugging this into $(1')$ and $(2')$, we get $\left( x' - y' \right) y' = 0$.

 Case 2.1: $y' = 0$.

 Thus, $z' = 0$. Plugging $y' = 0$ and $z' = 0$ into (1), we get a feasible solution $x = 11$, $y = 4$, $z = 4$.

 Case 2.2: $y' \neq 0$ and $x' = y'$.

 Thus, $x' = y' = z'$. Plugging these into (1), we get $\left( x, y, z \right) = \left( -10, -10, -10 \right)$ or $\left( 6, 6, 6 \right)$.

 Putting all cases together, $S = \left\{ 4, 11, -10, 6 \right\}$.Therefore, the sum of the squares of the elements of $S$ is
 \begin{align*} 4^2 + 11^2 + \left( -10 \right)^2 + 6^2 = \boxed{\textbf{(273) }}  . \end{align*}





\textbf{Solution 3 (Quadratic Formula and Vieta's Formulas)}

We index these equations as (1), (2), and (3), respectively.Using equation (1), we get $z = \frac{60 - xy}{4} = 15 - \frac{xy}{4}$We need to solve for x, so we plug this value of z into equation (3) to get:

 
 \[15x - \frac{x^2y}{4} - 4y = 60\]
 \[\frac{y}{4} * x^2 - 15x + (60 + 4y) = 0\]

 We use the quadratic formula to get possible values of x:

 
 \[x = \frac{15 \pm \sqrt{15^2 - 4(\frac{y}{4})(60 + 4y)}}{\frac{y}{2}}\]
 \[x = \frac{30 \pm 2\sqrt{225 - 60y + 4y^2}}{y}\]
 \[x = \frac{30 \pm 2\sqrt{(2y-15)^2}}{y}\]
 \[x = \frac{30 \pm 2(2y - 15)}{y}\]
 \[x = \frac{30 \pm (4y - 30)}{y}\]

 Here, we have two cases, (plus) and (minus)In the plus case, we have:
 \[x = \frac{(30 + 4y - 30)}{y}\]
 \[x = 4\]

 So, our first case gives us one value of x, which is 4.In the minus case, we have:
 \[x = \frac{30 - (4y - 30)}{y}\]
 \[x = \frac{60 - 4y}{y}\]
 \[x = \frac{60}{y} - 4\]

 For this case, we now have values of x in terms of y. Plugging this expression for x in equation (1), we get
 \[60 - 4y + 4z = 60\]
 \[4z = 4y\]
 \[z = y\]

 So we know that for this case, z = y. Using this information in equation (2), we get
 \[y^2 + \frac{240}{y} - 16 = 60\]
 \[y^2 + \frac{240}{y} - 76 = 0\]Multiplying both sides by y, we get a cubic expression:
 \[y^3 + 0y^2 - 76y + 240 = 0\]Here we just have to figure out the values of y that make this equation true. I used Vieta's Formulas to get a possible list, but you could also use the rational root theorem and synthetic division to find these. We call the three values of y that solve this equation: $y_{1},y_{2},y_{3}$Using Vieta's Formulas, you get these three expressions:
 \[y_{1} + y_{2} + y_{3} = 0\]
 \[y_{1} * y_{2} * y_{3} = -240\]
 \[(y_{1} * y_{2}) + (y_{2} * y_{3}) + (y_{1} * y_{3}) = -76\]

 In addition, we know that $y \vert 60$, because of our expression for x. Since the three values of y multiply to a negative number but also add to 0, we know that one value is negative and the other two are positive, and that the absolute value of the negative value is greater than both of the positive values.List of possible values for y are $\{-60,-30,-20,-15,-12,-10,-6,-5,-4,-3,-2,-1,1,2,3,4,5,6,10,12,15,20,30,60\}$From a list of these values, the only values that work are $y_1 = -10, y_2 = 6, y_3 = 4$ because
 \[-10 + 6 + 4 = 0\]
 \[-10 * 6 * 4 = -240\]
 \[(-10 * 6) + (-10 * 4) + (6 * 4) = -60 - 40 + 24 = -76\]

 Plugging in these values for y into our expression for x, we get:
 \[x = \frac{60}{-10} - 4 = -10\] 
 \[x = \frac{60}{6} - 4 = 6\] 
 \[x = \frac{60}{4} - 4 = 11\] 

 So, now we have accounted for both cases, and we have 4 values of x = $\{-10,4,6,11\}$Squaring all these terms we get: 100 + 16 + 36 + 121 = 273, so our answer is $\boxed{\textbf{(273)}}$




\textbf{Solution 4 (even more Vieta's)}

Since all three equations are in the form $\frac{K}{a} + 4a = 60$ where $K = xyz$, we can rearrange this to see that $x$, $y$, and $z$ all satisfy 

 
 \[ 4a^2 - 60a + K = 0. \]

 Let this quadratic have roots $a_1$ and $a_2$. Then, there are two cases to consider: two of $x$, $y$, $z$ are equal to $a_1$ and the third is equal to $a_2$, or all of $x$, $y$, $z$ are equal to $a_1$.

 Case 1: WLOG let $x = y = a_1$ and $z = a_2$.

 Then by Vieta's,

 
 \[x + z = 15 \hspace{1cm} (1)\]
 \[xz = \frac{x^2z}{4} \hspace{1cm} (2)\]

 which gives $x = y = 4$ and $z = 11$. But we can swap $x$, $y$, and $z$ however we like, so this also gives $x = 11$ as a solution. In total, this case yields $x = 4$ and $x = 11$ as possible values of $x$.

 Case 2: WLOG let $x = y = z = a_1$.

 Again, by Vieta's,

 
 \[x + a_2 = 15 \hspace{1cm} (1)\]
 \[xa_2 = \frac{x^3}{4} \hspace{1cm} (2)\]

 We can use $(2)$ to isolate $a_2$ in terms of $x$, then plug that into $(1)$ to get that $x^2 + 4x - 60 = 0$. This yields $x = -10$ and $x = 6$ as additional possible values of $x$.

 In all, $x$ can be any of $\{ 4, 11, -10, 6 \}$, so the requested answer is $4^2 + 11^2 + (-10)^2 + 6^2 = \boxed{273}.$




\textbf{Solution 5}

We index these equations as (1), (2), and (3) respectively (same as solution 2). There are two possible cases:

 Case 1: $x = \pm 4$

 In this case, we simply plug in $x = 4$ and $x = -4$. We note that $x=4$ is a valid case.

 Case 2: $x \neq \pm 4$

 In this case, using equation (3), we get $y = 15 - \frac{xz}{4}$. Plugging that into equation (1), we get $z = \frac{240-60x}{16-x^2}$. Plugging that expression back into the original expression for $y$ we obtain $y = 15 - \frac{60x-15x^2}{16-x^2}$. Now we plug these two expressions into equation (2):
 \[\left(15 - \frac{60x-15x^2}{16-x^2}\right)\left(\frac{240-60x}{16-x^2}\right) = 60-4x\]multiplying both sides by $(16-x^2)$ and factoring, we get:
 \[\left(15(4-x)(4+x)-60x+15x^2\right)\left(\frac{60(4-x)}{(4+x)(4-x)}\right) = 4(15 - x)(4 + x)(4 - x)\]which simplifies to:
 \[x^3 - 7x^2 - 104x + 660 = 0\]we note that $x = 6$ is a root. Factoring, we get the other roots, -10 and 11. 

 Our desired answer is the sum of the square of all these roots: 
 \[4^2 + 6^2 + (-10)^2 + 11^2 = \boxed{273}\]



	\item (\textit{AIME II 2023}) Let $S$ be the set of all positive rational numbers $r$ such that when the two numbers $r$ and $55r$ are written as fractions in lowest terms, the sum of the numerator and denominator of one fraction is the same as the sum of the numerator and denominator of the other fraction. The sum of all the elements of $S$ can be expressed in the form $\frac{p}{q},$ where $p$ and $q$ are relatively prime positive integers. Find $p+q.$

 \textbf{Answer: }719



\textbf{Solution 1}

Denote $r = \frac{a}{b}$, where $gcf\left( a, b \right) = 1$.We have $55 r = \frac{55a}{b}$.Suppose $gcf\left( 55, b \right) = 1$, then the sum of the numerator and the denominator of $55r$ is $55a + b$.This cannot be equal to the sum of the numerator and the denominator of $r$, $a + b$.Therefore, $gcf\left( 55, b \right) \neq 1$.

 Case 1: $b$ can be written as $5c$ with $gcf\left( c, 11 \right) = 1$.

 Thus, $55r = \frac{11a}{c}$.

 Because the sum of the numerator and the denominator of $r$ and $55r$ are the same,
 \[ a + 5c = 11a + c . \]

 Hence, $2c = 5 a$.

 Because $gcf\left( a, b \right) = 1$, $gcf\left( a, c \right) = 1$.Thus, $a = 2$ and $c = 5$.Therefore, $r = \frac{a}{5c} = \frac{2}{25}$.

 Case 2: $b$ can be written as $11c$ with $gcf\left( c, 5 \right) = 1$.

 Thus, $55r = \frac{5a}{c}$.

 Because the sum of the numerator and the denominator of $r$ and $55r$ are the same,
 \[ a + 11c = 5a + c . \]

 Hence, $2a = 5 c$.

 Because $gcf\left( a, b \right) = 1$, $gcf\left( a, c \right) = 1$.Thus, $a = 5$ and $c = 2$.Therefore, $r = \frac{a}{11c} = \frac{5}{22}$.

 Case 3: $b$ can be written as $55 c$.

 Thus, $55r = \frac{a}{c}$.

 Because the sum of the numerator and the denominator of $r$ and $55r$ are the same,
 \[ a + 55c = a + c . \]

 Hence, $c = 0$. This is infeasible.Thus, there is no solution in this case.

 Putting all cases together, $S = \left\{ \frac{2}{25}, \frac{5}{22} \right\}$.Therefore, the sum of all numbers in $S$ is
 \[ \frac{2}{25} + \frac{5}{22} = \frac{169}{550} . \]

 Therefore, the answer is $169 + 550 = \boxed{\textbf{(719) }}$.

 Minor latex edits by \href{https://artofproblemsolving.com/wiki/index.php?title=User:T3chn0b14d3&action=edit&redlink=1}{T3CHN0B14D3}

 



\textbf{Solution 2 (equation with GCDs)}

When we simplify a fraction $\frac{a}{b}$ over the integers, we divide the numerator and denominator by $\gcd(a,b).$ So we can clearly see that the fraction $r$ in lowest terms will be $\frac{(\frac{a}{\gcd(a,b)})}{(\frac{b}{\gcd(a,b)})}$and similarly, the fraction $55r$ in lowest terms will be $\frac{(\frac{55a}{\gcd(55a,b)})}{(\frac{b}{\gcd(55a,b)})}.$

 Now, from the problem's condition, we obtain the following equation: 
 \[\frac{a + b}{\gcd(a,b)} = \frac{55a + b}{\gcd(55a,b)}.\]

 Let $\gcd(a,b) = c,$ where $c$ is some positive integer. Then we can write $a$ as $cx$ and $b$ as $cy$ for positive integers $x$ and $y$ such that $\gcd(x,y) = 1.$ Substituting in $cx$ for $a$ and $cy$ for $b$ we get 
 \[\frac{cx + cy}{c} = \frac{55cx + cy}{\gcd(55cx, cy)}\]

 which simplifies to 

 
 \[x + y = \frac{55x + y}{\gcd(55x,y)}\]

 Now, note that since $x$ and $y$ share no common factors, any common factors shared between $55x$ and $y$ must be factors of $55,$ so from here we do casework on the possible values of $\gcd(55,y).$ We know that $\gcd(55,y) \in \{1, 5, 11, 55\},$ so we only have to look at four different cases. 

 Case 1: $\gcd(55,y) = 1:$ This implies 
 \[x + y = 55x + y\] which would make one of the variables $0.$ This case doesn't work. 

 Case 2: $\gcd(55,y) = 5:$ 
 \[5x + 5y = 55x + y \implies 50x = 4y \implies x = \frac{2}{25}y.\] Since $a = cx$ and $b = cy,$ the value of $r$ is just $\frac{a}{b} = \frac{x}{y},$ so this case gives $\frac{2}{25}$ as a possible value for $r$. 

 Case 3: $\gcd(55,y) = 11:$ 
 \[11x + 11y = 55x + y \implies 44x = 10y \implies 22x = 5y \implies r = \frac{5}{22}.\] 

 Case 4: $\gcd(55,y) = 55:$ 
 \[55x + 55y = 55x + y\] This also makes one of the variables $0$ which we can't have. 

 So the answer is $\frac{2}{25} + \frac{5}{22} = \frac{169}{550}.$
$169 + 550 = \boxed{719}.$




\textbf{Solution 3(Easiest)}

Write $r = \frac{p}{q}, \gcd(p, q) = 1 \implies 55r = \frac{55p}{q}$. 

 Case 1: $\gcd(q, 55) = 1$

 Implies $p + q = 55p + q \implies p = 0$ a contradiction.

 Case 2: $\gcd(q, 11) = 1, \gcd(q, 5) = d$ where $d > 1$.

 Implies $q = 5k$ where $\gcd(k, 11) = 1$. This means $55r = \frac{11p}{k} \implies p + q = 11p + k \implies p + 5k = 11p + k \implies 4k = 10p \implies 2k = 5p$ since $\gcd(p, q) = 1 \implies \gcd(p, k) = 1 \implies p = 2, k = 5 \implies r = \frac{p}{5k} = \frac{2}{25}$.

 Case 3: $\gcd(q, 5) = 1, \gcd(q, 11) = d$ where $d > 1$.

 Implies $q = 11k$ where $\gcd(k, 11) = 1$. This means $55r = \frac{5p}{k} \implies p + q = p + 11k = 5p + k \implies 10k = 4p \implies 5k = 2p$. Since $\gcd(p, q) = 1 \implies \gcd(p, k) = 1 \implies k = 2, p = 5 \implies r = \frac{p}{11k} = \frac{5}{22}$.

 Case 4: $\gcd(q, 55) = d$ where $d > 1$.

 Implies $q = 55k$ where $\gcd(k, 55) = 1$. This means $55r = \frac{p}{k} \implies p + q = p + k \implies q = 55k = k \implies k = 0$ a contradiction.

 So $\frac{2}{25} + \frac{5}{22} = \frac{169}{550} \implies \boxed{719}$.

 \textbf{Note} This problem mainly comes down to noticing that $55r$ has to be simplifiable such that the numerator and denominator both change, so they potentially equal their original sum. Then you proceed with casework just as Solution 1.

 


	\item (\textit{AIME II 2023}) Consider the L-shaped region formed by three unit squares joined at their sides, as shown below. Two points $A$ and $B$ are chosen independently and uniformly at random from inside the region. The probability that the midpoint of $\overline{AB}$ also lies inside this L-shaped region can be expressed as $\frac{m}{n},$ where $m$ and $n$ are relatively prime positive integers. Find $m+n.$

                    \begin{figure}[H]
                        \centering
                        \includegraphics[width=0.5\textwidth]{Problem/AIME_II_2023_6_Problem_Figure_1.png}
                    \end{figure}
                    

 \textbf{Answer: }035



\textbf{Solution 1}

We proceed by calculating the complement. 

 Note that the only configuration of the 2 points that makes the midpoint outside of the L shape is one point in the top square, and one in the right square. This occurs with $\frac{2}{3} \cdot \frac{1}{3}$ probability. 

 Let the topmost coordinate have value of: $(x_1,y_1+1)$, and rightmost value of: $(x_2+1,y_2)$.

 The midpoint of them is thus: 
 \[\left(\frac{x_1+x_2+1}{2}, \frac{y_1+y_2+1}{2} \right)\]

 It is clear that $x_1, x_2, y_1, y_2$ are all between 0 and 1. For the midpoint to be outside the L-shape, both coordinates must be greater than 1, thus: 
 \[\frac{x_1+x_2+1}{2}>1 \implies x_1+x_2>1\]

 By symmetry this has probability $\frac{1}{2}$. Also by symmetry, the probability the y-coordinate works as well is $\frac{1}{2}$. Thus the probability that the midpoint is outside the L-shape is: 
 \[\frac{2}{3} \cdot \frac{1}{3} \cdot \frac{1}{2} \cdot \frac{1}{2} = \frac{1}{18}\]

 We want the probability that the point is inside the L-shape however, which is $1-\frac{1}{18}=\frac{17}{18}$, yielding the answer of $17+18=\boxed{35}$~SAHANWIJETUNGA

 



\textbf{Solution 2(Calculus)}

We assume each box has side length 1.We index the upper left box, the bottom left box, the bottom right box as II, III, IV, respectively.We index the missing upper right box as I.We put the graph to a coordinate system by putting the intersecting point of four foxes at the origin, the positive direction of the $x$axis at the intersecting line of boxes I and IV, and the positive direction of the $y$-axis at the intersecting line of boxes I and II.We denote by $M$ the midpoint of $AB$.

 Therefore,
 \begin{align*} \Bbb P \left( M \in II \cup III \cup IV \right) & = 1 - \Bbb P \left( M \in I \right) . \end{align*}


 We observe that a necessary for $M \in I$ is either $A \in II$ and $B \in IV$, or $A \in IV$ and $B \in II$.In addition, by symmetry,
 \[ \Bbb P \left( M \in I | A \in II, B \in IV \right) = \Bbb P \left( M \in I | A \in IV, B \in II \right) . \]

 Thus,
 \begin{align*} \Bbb P \left( M \in I \right) & = 2 \Bbb P \left( M \in I | A \in II, B \in IV \right) \Bbb P \left( A \in II, B \in IV \right) \\ & = 2 \Bbb P \left( M \in I | A \in II, B \in IV \right) \Bbb P \left( A \in II \right) \Bbb P \left( B \in IV \right) \\ & = 2 \Bbb P \left( M \in I | A \in II, B \in IV \right) \cdot \frac{1}{3} \cdot \frac{1}{3} \\ & = \frac{2}{9} \Bbb P \left( M \in I | A \in II, B \in IV \right) \\ & = \frac{2}{9} \Bbb P \left( x_M > 0, y_M > 0 | x_A < 0, y_A > 0, x_B > 0, y_B < 0 \right) \\ & = \frac{2}{9} \Bbb P \left( x_M > 0 | x_A < 0, x_B > 0 \right) \cdot \Bbb P \left( y_M > 0 | y_A > 0, y_B < 0 \right) \\ & = \frac{2}{9} \Bbb P \left( x_M > 0 | x_A < 0, x_B > 0 \right)^2 \\ & = \frac{2}{9} \Bbb P \left( x_A + x_B > 0 | x_A < 0, x_B > 0 \right)^2 \\ & = \frac{2}{9} \left( \int_{x_A = -1}^0 \int_{x_B = 0}^1 \mathbf 1 \left\{ x_A + x_B > 0 \right\} dx_B dx_A \right)^2 \\ & = \frac{2}{9} \left( \int_{x_A = -1}^0 \int_{x_B = - x_A}^1 dx_B dx_A \right)^2 \\ & = \frac{2}{9} \left( \int_{x_A = -1}^0 \left( 1 + x_A \right) dx_A \right)^2 \\ & = \frac{2}{9} \left( \frac{1}{2}  \right)^2 \\ & = \frac{1}{18} . \end{align*}
The second equality follows from the condition that the positions of $A$ and $B$ are independent.The sixth equality follows from the condition that for each point of $A$ and $B$, the $x$ and $y$ coordinate are independent.

 Therefore,
 \begin{align*} \Bbb P \left( M \in II \cup III \cup IV \right) & = 1 - \frac{1}{18} \\ & = \frac{17}{18} . \end{align*}


 Therefore, the answer is $17 + 18 = \boxed{\textbf{(035) }}$.

 



\textbf{Solution 3 (Geometry)}


                    \begin{figure}[H]
                        \centering
                        \includegraphics[width=0.5\textwidth]{Solution/AIME_II_2023_6_Solution_Figure_1.png}
                    \end{figure}
                    The only configuration of two points that makes the midpoint outside of the $L-$shape is one point in the top square, and one in the right square with probability $\frac{2}{9}$ (see Solution $1.)$ 

 We use the coordinate system shown in diagram. Let the arbitrary point $A(x, -y), x \in [0,1], y \in [0,1]$ be in right square. It is clear that iff point $B$ lies in yellow rectangle with sides $x$ and $1 – y,$ then the midpoint $AB$ lies outside of the $L.$

 Probability of this is $p = x(1 – y).$

 Consider the points $A_1(1-y,-x), A_2(1-x,y-1), A_3(y,x-1).$

 Similarly we find $p_1 = (1-y)(1-x), p_2= (1-x)y,  p_3 = xy.$ 

 The probability that the midpoint of one of the segments $A_1B, A_2B, A_3B$ and $AB$ is outside of the $L-$shape is 
 \[\frac {p_1+p_2+p_3+p}{4} = \frac {1}{4}.\]If point $A$ is in the right square, point $B$ in the top square, the probability that the midpoint of $\overline{AB}$ lies outside $L-$shape is $\frac {1}{4} \implies$

 If points $A$ and $B$ are chosen independently and uniformly at random in $L-$shape, then the probability that the midpoint of $\overline{AB}$ lies outside $L-$shape is $\frac {1}{4} \cdot \frac {2}{9} = \frac {1}{18}.$Therefore the probability that the point is inside the $L-$shape is $1-\frac{1}{18}=\frac{17}{18} \implies \boxed{35}.$

 \textbf{vladimir.shelomovskii@gmail.com, vvsss}

 



\textbf{Solution 4 (Coordinates, similar to Solution 1)}

Consider this diagram:

 
                    \begin{figure}[H]
                        \centering
                        \includegraphics[width=0.5\textwidth]{Solution/AIME_II_2023_6_Solution_Figure_2.png}
                    \end{figure}
                    

 First, the one of points must be in the uppermost box and the other in the rightmost box. This happens with probability $\frac{2}{3}*\frac{1}{3}=\frac{2}{9}$.

 We need the midpoints of the $x$ coordinates to be greater than $1$ but less than $2.$ 

 We need the midpoints of the $y$ coordinates to be greater than $1$ but less than $2.$ 

 Thus, we set up:
 \[1<\frac{x_1+x_2}{2}<2\]
 \[1<\frac{y_1+y_2}{2}<2\]

 Using geometric probability, as shown below, we get that the probability of happening is $\frac{1}{2}$. This is for both. So we have probability $\frac{1}{2}*\frac{1}{2}=\frac{1}{4}$.

 
                    \begin{figure}[H]
                        \centering
                        \includegraphics[width=0.5\textwidth]{Solution/AIME_II_2023_6_Solution_Figure_3.png}
                    \end{figure}
                    

 Thus, the fraction is $1-(\frac{2}{9}*\frac{1}{4})= 1-\frac{1}{18}=\frac{17}{18}$.

 The answer is $17+18=35$

 mathboy282

 Minor edit by whatsoever281

 



\textbf{Solution 5}

We casework on the squares the two points are in. Unless the two squares are on opposite corners, which happens with probability $\frac{2}{9}$. If the points are on opposite corners, the points must be on the upper half of the splitting diagonal for the midpoint to not be in the square. This occurs with probability $\frac{1}{4}$. We take the complement and get final answer $(\frac{2}{9})(\frac{3}{4}) + (\frac{7}{9}) = \frac{17}{18}$.

 CORRECTION: This solution is wrong. We can have one point above the splitting diagonal and one point below the splitting diagonal and still have the midpoint not be in the square.

 


	\item (\textit{AIME II 2023}) Each vertex of a regular dodecagon ($12$-gon) is to be colored either red or blue, and thus there are $2^{12}$ possible colorings. Find the number of these colorings with the property that no four vertices colored the same color are the four vertices of a rectangle.


 
 \textbf{Answer: }928



\textbf{Solution 1}

Note that the condition is equivalent to stating that there are no 2 pairs of oppositely spaced vertices with the same color.

 Case 1: There are no pairs. This yields $2$ options for each vertices 1-6, and the remaining vertices 7-12 are set, yielding $2^6=64$ cases.

 Case 2: There is one pair. Again start with 2 options for each vertex in 1-6, but now multiply by 6 since there are 6 possibilities for which pair can have the same color assigned instead of the opposite. Thus, the cases are: $2^6*6=384$

 case 3: There are two pairs, but oppositely colored. Start with $2^6$ for assigning 1-6, then multiply by 6C2=15 for assigning which have repeated colors. Divide by 2 due to half the cases having the same colored opposites. $2^6*15/2=480$

 It is apparent that no other cases exist, as more pairs would force there to be 2 pairs of same colored oppositely spaced vertices with the same color. Thus, the answer is: $64+384+480=\boxed{928}$




\textbf{Solution 2}

First, we identify the geometric condition for the sentence ``four vertices colored thesame color are the four vertices of a rectangle\textit{.}Consider any four vertices on the dodecagon, $A$, $B$, $C$, $D$.Denote by $O$ the center of the dodecagon.Because $OA = OB = OC$, $\angle OAB = \angle OBA$ and $\angle OBC = \angle OCB$.

 Thus,
 \begin{align*} \angle ABC & = \angle OBA + \angle OBC \\ & = \frac{\angle OBA + \angle OAB}{2} + \frac{\angle OBC + \angle OCB}{2} \\ & = \frac{180^\circ - \angle AOB}{2} +  \frac{180^\circ - \angle COB}{2} \\ & = 180^\circ - \frac{\angle AOB + \angle COB}{2} \\ & = 180^\circ - \frac{\angle AOC}{2} . \end{align*}


 Hence, $\angle ABC = 90^\circ$ if and only if $\angle AOC = 180^\circ$.Similarly, $\angle ADC = 90^\circ$ if and only if $\angle AOC = 180^\circ$, and $\angle BCD = 90^\circ$ (or $\angle DAB = 90^\circ$) if and only if $\angle BOD = 180^\circ$.

 Therefore, $ABCD$ is a rectangle if and only if two diagonals both pass through $O$.

 Now, we categorize 12 vertices into 6 groups. Each group contains 2 diagonal vertices.Next, we compute the number of coloring configurations such that the above same-color rectangles do not exist.

 Case 1: Two vertices in each group has distinct colors.

 For each group, we only need to determine which vertex is red. The other one must be blue.Therefore, the number of configurations in this case is $2^6$.

 Case 2: There is one group who vertices have the same color. All other groups are with vertices that have distinct colors.

 We construct such an instance in the following steps.

 Step 1: We determine which group has two vertices that have the same color.

 The number of ways is 6.

 Step 2: For the selected group, we choose a color for its two vertices.

 The number of ways is 2.

 Step 3: For each unselected group, we determine which vertex is red.

 The number of ways is $2^5$.

 Following from the rule of product, the total number of configurations in this case is $6 \cdot 2 \cdot 2^5 = 6 \cdot 2^6$.

 Case 3: One group has two red vertices, one group has two blue vertices, and each of the other four groups has vertices with distinct colors.

 We construct such an instance in the following steps.

 Step 1: We determine which group has two vertices that have both red color.

 The number of ways is 6.

 Step 2: We determine which group has two vertices that have both blue color.

 The number of ways is 5.

 Step 3: For each unselected group, we determine which vertex is red.

 The number of ways is $2^4$.

 Following from the rule of product, the total number of configurations in this case is $6 \cdot 5 \cdot 2^4 = 30 \cdot 2^4$.

 Putting all cases together, the total number of feasible configurations is
 \[ 2^6 + 6 \cdot 2^6 + 30 \cdot 2^4 = \boxed{\textbf{(928) }}. \]




\textbf{Solution 3 (Similar to Solution 1, just reworded a bit)}

Note that a rectangle is formed if and only if it's diagonals pass through the center of the dodecagon and the diagonal's endpoints are the same color. 

 Consider the $6$ diagonals that pass through the center. A rectangle is formed if the endpoints of some pair of them are all the same color. We can now perform casework on the diagonals:

 Case $1:$ The endpoints of all the diagonals are all different (e.g. for all diagonals, the endpoints are red and blue in some order).

 The amount of ways to do this is $2^6=64,$ since there are $2$ ways to order which vertice is red and which is blue for each diagonal. 

 Case $2:$ There exists one diagonal such that the endpoints are the same color.

 There are $6$ ways to choose this diagonal, $2$ ways to choose which color it is, and $2^5$ ways to color the rest of the diagonals. Therefore, the total for this case is $6 \cdot 2 \cdot 2^5 = 384.$

 Case $3:$ There exists two diagonals such that the endpoints are the same color, but they are different colors

 Then, there are $6 \cdot 5=30$ ways to choose these two diagonals, and $2^4=16$ ways to color the rest, so the total for this case is $30 \cdot 16 = 480.$ 

 Summing all of the cases, we have $64+384+480=\boxed{928}.$




\textbf{Solution 4 \approx Solution 1}


 \[\text{First, we notice that a rectangle is made from two pairs of vertices 1/2 turn away from each other.}\]

                    \begin{figure}[H]
                        \centering
                        \includegraphics[width=0.5\textwidth]{Solution/AIME_II_2023_7_Solution_Figure_1.png}
                    \end{figure}
                    
 \[\textit{Note: The image is }\frac{\textit{280}}{\textit{841}}\approx\frac{\textit{1}}{\textit{3}}\textit{ size.}\]

 \[\text{For there to be no rectangles, there can be at most one same-colored pair for each color and the rest need to have one red and blue.}\]

 \[\textit{\underline{Case 1: \textbf{No pairs}}}\]

 \[\text{Each pair has two ways to color: One red or the other red. Thus, there are }2^6=\boxed{64}\text{ ways in this case.}\]

 \[\textit{\underline{Case 2: \textbf{One red pair}}}\]

 \[\text{The red pair has }\binom{6}{1}\text{ positions. All the rest still have two ways. Therefore, there are }\binom{6}{1}\cdot 2^5=\frac{6}{1}\dot 2^5=6\cdot 32=\boxed{192} \text{ ways in this case.}\]

 \[\textit{\underline{Case 3: \textbf{One blue pair}}}\]

 \[\text{This is the same as the one red pair case so there are still }\binom{6}{1}\cdot 2^5=6\cdot 2^5=6\cdot 32=\boxed{192}\text{ ways.}\]

 \[\textit{\underline{Case 4: \textbf{One pair of each color}}}\]

 \[\text{The red pair has }\binom{6}{1}\text{ positions. The blue pair has }\binom{5}{1}\text{ positions. All the rest still have two ways. Therefore, there are }\binom{6}{1}\cdot\binom{5}{1}\cdot 2^4=\frac{6\cdot 5=30}{1\cdot 1=1}\cdot 2^4=30\cdot 16=\boxed{480}\text{ ways in this case.}\]

 \[\textit{\underline{\textbf{Solution}}}\]

 \[\text{In total, there are }64+192+192+480=\boxed{928}\text{ways.}\]



	\item (\textit{AIME II 2023}) Let $\omega = \cos\frac{2\pi}{7} + i \cdot \sin\frac{2\pi}{7},$ where $i = \sqrt{-1}.$ Find the value of the product 
 \[\prod_{k=0}^6 \left(\omega^{3k} + \omega^k + 1\right).\]

 \textbf{Answer: }024



\textbf{Solution 1}

For any $k\in \mathbb{Z}, 7\nmid k$, we have,
 \begin{align*} & \left( \omega^{3k} + \omega^k + 1 \right) \left( \omega^{3\left( 7 - k \right)} + \omega^{\left( 7 - k \right)} + 1 \right) \\ & = \omega^{3 \cdot 7} + \omega^{2k + 7} + \omega^{3k} + \omega^{-2k + 3 \cdot 7} + \omega^7 + \omega^k + \omega^{3\left( 7 - k \right)} + \omega^{\left( 7 - k \right)} + 1 \\ & = 1 + \omega^{2k} + \omega^{3k} + \omega^{-2k} + 1 + \omega^k + \omega^{-3k} + \omega^{-k} + 1 \\ & = 2 + \omega^{-3k} \sum_{j=0}^6 \omega^{j k} \\ & = 2 + \omega^{-3k} \frac{1 - \omega^{7 k}}{1 - \omega^k} \\ & = 2 . \end{align*}
The second and the fifth equalities follow from the property that $\omega^7 = 1$.

 Therefore,
 \begin{align*} \Pi_{k=0}^6 \left( \omega^{3k} + \omega^k + 1 \right) & = \left( \omega^{3 \cdot 0} + \omega^0 + 1 \right) \Pi_{k=1}^3 \left( \omega^{3k} + \omega^k + 1 \right) \left( \omega^{3\left( 7 - k \right)} + \omega^{\left( 7 - k \right)} + 1 \right) \\ & = 3 \cdot 2^3 \\ & = \boxed{\textbf{024}}. \end{align*}

$k$ cannot be a multiple of $7$, otherwise the first equation would equal $9$ instead of $2$ ~inaccessibles

 



\textbf{Solution 2 (Moduli)}

Because the answer must be a positive integer, it is just equal to the modulus of the product. Define $z_n = \left(\textrm{cis }\frac{2n\pi}{7}\right)^3 + \textrm{cis }\frac{2n\pi}{7} + 1$.

 Then, our product is equal to

 
 \[|z_0||z_1||z_2||z_3||z_4||z_5||z_6|.\]
$z_0 = 3$, and we may observe that $z_x$ and $z_{7-x}$ are conjugates for any $x$, meaning that their magnitudes are the same. Thus, our product is

 
 \[3|z_1|^2|z_2|^2|z_3|^2\]
 \[= 3\left((\cos \frac{6\pi}{7} + \cos \frac{2\pi}{7} + 1)^2 + (\sin \frac{6\pi}{7} + \sin \frac{2\pi}{7})^2\right) \left((\cos \frac{12\pi}{7} + \cos \frac{4\pi}{7} + 1)^2 + (\sin \frac{12\pi}{7} + \sin \frac{4\pi}{7})^2\right) \left((\cos \frac{4\pi}{7} + \cos \frac{6\pi}{7} + 1)^2 + (\sin \frac{4\pi}{7} + \sin \frac{6\pi}{7})^2\right)\]

 Let us simplify the first term. Expanding, we obtain

 
 \[\cos^2 \frac{6\pi}{7} + \cos^2 \frac{2\pi}{7} + 1 + 2\cos \frac{6\pi}{7}\cos \frac{2\pi}{7} + 2\cos \frac{6\pi}{7} + 2\cos \frac{2\pi}{7} + \sin^2 \frac{6\pi}{7} + \sin^2 \frac{2\pi}{7} + 2\sin \frac{6\pi}{7}\sin \frac{2\pi}{7}.\]

 Rearranging and cancelling, we obtain

 
 \[3 + 2\cos \frac{6\pi}{7} + 2\cos \frac{2\pi}{7} + 2\cos \frac{6\pi}{7}\cos \frac{2\pi}{7} + 2\sin \frac{6\pi}{7}\sin \frac{2\pi}{7}.\]

 By the cosine subtraction formula, we have $2\cos \frac{6\pi}{7}\cos \frac{2\pi}{7} + 2\sin \frac{6\pi}{7}\sin \frac{2\pi}{7} = \cos \frac{6\pi - 2\pi}{7} = \cos \frac{4\pi}{7}$.

 Thus, the first term is equivalent to 

 
 \[3 + 2(\cos \frac{2\pi}{7} + \cos \frac{4\pi}{7} + \cos \frac{6\pi}{7}).\]

 Similarly, the second and third terms are, respectively, 

 
 \[3 + 2(\cos \frac{4\pi}{7} + \cos \frac{8\pi}{7} + \cos \frac{12\pi}{7}),\textrm{ and}\]
 \[3 + 2(\cos \frac{6\pi}{7} + \cos \frac{12\pi}{7} + \cos \frac{4\pi}{7}).\]

 Next, we have $\cos \frac{2\pi}{7} + \cos \frac{4\pi}{7} + \cos \frac{6\pi}{7} = -\frac{1}{2}$. This is because

 
 \[\cos \frac{2\pi}{7} + \cos \frac{4\pi}{7} + \cos \frac{6\pi}{7} = \frac{1}{2}(\textrm{cis }\frac{2\pi}{7} + \textrm{cis }\frac{4\pi}{7} + \textrm{cis }\frac{6\pi}{7} + \textrm{cis }\frac{8\pi}{7} + \textrm{cis }\frac{10\pi}{7} + \textrm{cis }\frac{12\pi}{7})\]

 \[= \frac{1}{2}(-1)\]
 \[= -\frac{1}{2}.\]

 Therefore, the first term is simply $2$. We have $\cos x = \cos 2\pi - x$, so therefore the second and third terms can both also be simplified to $3 + 2(\cos \frac{2\pi}{7} + \cos \frac{4\pi}{7} + \cos \frac{6\pi}{7}) = 2$. Thus, our answer is simply

 
 \[3 \cdot 2 \cdot 2 \cdot 2\]
 \[= \boxed{\mathbf{024}}.\]




\textbf{Solution 3 (Inspecting the exponents of powers of )}

We write out the product in terms of $\omega$:
 \[\prod_{k=0}^6 \left(\omega^{3k} + \omega^k + 1\right)=3(\omega^3+\omega+1)(\omega^6+\omega^2+1)(\omega^9+\omega^3+1)(\omega^{12}+\omega^4+1)(\omega^{15}+\omega^5+1)(\omega^{18}+\omega^6+1).\]

 Grouping the terms in the following way exploits the fact that $\omega^{7k}=1$ for an integer $k$, when multiplying out two adjacent products from left to right:

 
 \[\frac{1}{3} \prod_{k=0}^6 \left(\omega^{3k} + \omega^k + 1\right)=(\omega^3+\omega+1)(\omega^{18}+\omega^6+1)(\omega^6+\omega^2+1)(\omega^{15}+\omega^5+1)(\omega^9+\omega^3+1)(\omega^{12}+\omega^4+1).\]

 
When multiplying two numbers with like bases, we add the exponents. We can now rewrite the exponents of each product (two at a time, where $1$ is treated as the identity) as a series of arrays: 

 
 \[\textbf{(A)}\begin{bmatrix} 3&1 &0 \\ 18&6&0\\ \end{bmatrix}\]

 \[\textbf{(B)}\begin{bmatrix} 6&2 &0 \\ 15&5&0\\ \end{bmatrix}\]

 \[\textbf{(C)}\begin{bmatrix} 9&3 &0 \\ 12&4&0\\ \end{bmatrix}.\]

 
Note that $\omega=e^{\frac{2\pi i}{7}}$. When raising $\omega$ to a power, the numerator of the fraction is $2$ times whatever power $\omega$ is raised to, multiplied by $\pi i$. Since the period of $\omega$ is $2\pi,$ we multiply each array by $2$ then reduce each entry $\mod{14},$ as each entry in an array represents an exponent which $\omega$ is raised to.

 

 \[\textbf{(A)}\begin{bmatrix} 6&2 &0 \\ 8&12&0\\ \end{bmatrix}\]

 \[\textbf{(B)}\begin{bmatrix} 12&4 &0 \\ 2&10&0\\ \end{bmatrix}\]

 \[\textbf{(C)}\begin{bmatrix} 4&6 &0 \\ 10&8&0\\ \end{bmatrix}.\]

 To obtain the correct exponents, we seperately add each element of the lower row to one element of the top row. 

 Therefore (after reducing $\mod 14$ again), we get the following sets: 

 
 \[\textbf{(A)}\ \{0, 4, 6, 10, 0, 2, 8, 12, 0\}\]
 \[\textbf{(B)}\ \{0, 8, 12, 6, 0, 4, 2, 10, 0\}\]
 \[\textbf{(C)}\ \{0, 12, 4, 2, 0, 8, 10, 8, 0\}.\]

 Raising $\omega$ to the power of each element in every set then multiplying over $\textbf{(A)}, \textbf{(B)},$ and $\textbf{(C)}$ yields 

 
 \[\frac{1}{3} \prod_{k=0}^6 \left(\omega^{3k} + \omega^k + 1\right)=\left(\sum_{a\in \textbf{(A)}} \omega^a\right)\left(\sum_{b\in \textbf{(B)}} \omega^b\right)\left(\sum_{c\in \textbf{(C)}} \omega^c\right)\]

 \[=\left(\sum_{a\in \textbf{(A)}} \omega^a\right)^3\]

 \[=\left(\omega^0+\omega^4+\omega^6+\omega^{10}+\omega^0+\omega^2+\omega^8+\omega^{12}+\omega^0\right)^3\]

 \[=\left(3+\omega^2+\omega^4+\omega^6+\omega^8+\omega^{10}+\omega^{12}\right)^3,\]as these sets are all identical. 

 Summing as a geometric series, 

 
 \[\frac{1}{3} \prod_{k=0}^6 \left(\omega^{3k} + \omega^k + 1\right)=\left(3+\frac{\omega^2(\omega^{12}-1)}{\omega^2-1}\right)^3\]

 \[=\left(3+\frac{\omega^{14}-\omega^2}{\omega^2-1}\right)^3\]

 \[=\left(3+\frac{1-\omega^2}{\omega^2-1}\right)^3\]

 \[=(3-1)^3=8.\]

 Therefore, 

 
 \[\frac{1}{3} \prod_{k=0}^6 \left(\omega^{3k} + \omega^k + 1\right)=8,\]and 
 \[\prod_{k=0}^6 \left(\omega^{3k} + \omega^k + 1\right)=3\cdot8=\boxed{\textbf{(024)}}.\]

 -Benedict T (countmath1)

 


 



\textbf{Solution 4}

The product can be factored into $(1-r)(1-s)(1-t)(w-r)(w-s)(w-t)(w^2-r)(w^2-s)(w^2-t)....(w^6-r)(w^6-s)(w^6-t)$, (treat $\omega^k$ as $x$, so $x^3+x+1=0$ then k varies from 0 to 6.

 
where $r,s,t$ are the roots of the polynomial $x^3+x+1=0$. 

 
This is then $-(r^7-1)(s^7-1)(t^7-1)$ because $(r^7-1)$ and $(r-1)(r-w)(r-w^2)...(r-w^6)$ share the same roots and leading coefficient so are thus the same polynomial. 

 
To find $-(r^7-1)(s^7-1)(t^7-1)$,

 
Notice that $(r^7-1)=(r-1)(r^6+r^5+r^4+r^3+r^2+r+1)$. Since r satisfies $x^3+x+1=0$, $r^6+r^4+r^3=0$ 

 
Substituting, you are left with $r^5+r^2+r+1$. This is $r^2(r^3+1)+r+1$, and after repeatedly substituting $r^3+r+1=0$ you are left with $-2r^3$. 

 
So now the problem is reduced to finding $-(r-1)(s-1)(t-1)(-2r^3)(-2s^3)(-2t^3)=8(rst)^3(r-1)(s-1)(t-1)$, and vietas gives you the result of $\boxed{24}$ -resources

 



\textbf{Solution 5}

The product can be written as $(\omega^0+\omega^0+1)(\omega^3+\omega+1)(\omega^6+\omega^2+1)(\omega^9+\omega^3+1)\\(\omega^{12}+\omega^4+1)(\omega^{15}+\omega^5+1)(\omega^{18}+\omega^6+1).$ 

 The key here then is noticing that cis$2\pi=1.$ 

 From this we realize $\omega^{7n+a}=\omega^a$ where $a$ is a residue$\pmod 7.$

 Our expression then simplifies to $(3)(\omega^3+\omega+1)(\omega^6+\omega^2+1)(\omega^2+\omega^3+1)\\(\omega^5+\omega^4+1)(\omega+\omega^5+1)(\omega^4+\omega^6+1).$

 The key then becomes multiplying the 2nd and 7th, 3rd and 6th, 4th and 5th terms because they would result in a couple of $\omega^7$ which will cancel out.

 Multiplying them you obtain $\omega^6+\omega^5+\omega^4+\omega^3+\omega^2+\omega+3$. Recognizing this as $\frac{\omega^7-1}{\omega-1}+2.$ For terms $2*7, 3*6, 4*5.$
$\omega^7-1=0$ (cis$2\pi=1.$) We are left with $(3)(2)(2)(2)=\boxed{024}.$




\textbf{Solution 6 (Simple with proof)}

Let \( x = \omega^k \). Let \( \alpha_1 \), \( \alpha_2 \), and \( \alpha_3 \) be the roots of the polynomial \( x^3 + x + 1 \). Thus, the product simplifies down to
 \[\prod_{k=0}^6 (x - \alpha_1)(x - \alpha_2)(x - \alpha_3)\]Return \( x \) as the value \( \omega^k \). It becomes evident that every \( \omega^k \) is also a 7th root of unity. Thus, because each product cycles between the seven roots of unity, one can rewrite the entirety of the product as \( -(\alpha_1^7 - 1)(\alpha_2^7 - 1)(\alpha_3^7 - 1) \). Through Vieta's, \( \alpha_1 \alpha_2 \alpha_3 = -1 \), and through \href{https://artofproblemsolving.com/wiki/index.php?title=Newton\%27s\_Sums}{Newton's Sums}, the answer is $\boxed{024}$.

 



\textbf{Solution 6 Proof}

The most likely confusion point is: how do we get from simplifying the product into $(x-\alpha_1)(x-\alpha_2)(x-\alpha_3)$ into immediately concluding that the product is essentially $-(\alpha_1^7-1)(\alpha_2^7-1)(\alpha_3^7-1)$?

 We can prove this quite easily.

 Let $x$ be an $n$th root of unity, so $x^n = 1$.

 Factor $a^n - 1$ as $(a-1)(a-x)(a-x^2) \cdots (a-x^{n-1}) \square$.

 For each term in the product, factor out $-1$ in each product to obtain $(-1)^{n}\cdot (1-a)(x-a)(x^2-a)(x^3-a)(...)(x^{n-1}-a)$.

 Because $n=7$, we have the term $(-1)^{7} = -1$.

 Now, note that each $a$ term is a root of unity! This is because $a=1$, and that must correspond to some $x$ power. However, $x$ is also a root of unity! Thus, $a$ must also be a root of unity. Then, because $n=7$, our final rewritten form is $-(a^7-1)$.

 So, for our case

 


	\item (\textit{AIME II 2023}) Circles $\omega_1$ and $\omega_2$ intersect at two points $P$ and $Q,$ and their common tangent line closer to $P$ intersects $\omega_1$ and $\omega_2$ at points $A$ and $B,$ respectively. The line parallel to $AB$ that passes through $P$ intersects $\omega_1$ and $\omega_2$ for the second time at points $X$ and $Y,$ respectively. Suppose $PX=10,$ $PY=14,$ and $PQ=5.$ Then the area of trapezoid $XABY$ is $m\sqrt{n},$ where $m$ and $n$ are positive integers and $n$ is not divisible by the square of any prime. Find $m+n.$

 \textbf{Answer: }033



\textbf{Solution 1}

Denote by $O_1$ and $O_2$ the centers of $\omega_1$ and $\omega_2$, respectively.Let $XY$ and $AO_1$ intersect at point $C$.Let $XY$ and $BO_2$ intersect at point $D$.

 Because $AB$ is tangent to circle $\omega_1$, $O_1 A \perp AB$.Because $XY \parallel AB$, $O_1 A \perp XY$.Because $X$ and $P$ are on $\omega_1$, $O_1A$ is the perpendicular bisector of $XP$.Thus, $PC = \frac{PX}{2} = 5$.

 Analogously, we can show that $PD = \frac{PY}{2} = 7$.

 Thus, $CD = CP + PD = 12$.Because $O_1 A \perp CD$, $O_1 A \perp AB$, $O_2 B \perp CD$, $O_2 B \perp AB$, $ABDC$ is a rectangle. Hence, $AB = CD = 12$.

 Let $QP$ and $AB$ meet at point $M$.Thus, $M$ is the midpoint of $AB$.Thus, $AM = \frac{AB}{2} = 6$. This is the case because $PQ$ is the radical axis of the two circles, and the powers with respect to each circle must be equal.

 In $\omega_1$, for the tangent $MA$ and the secant $MPQ$, following from the power of a point, we have $MA^2 = MP \cdot MQ$.By solving this equation, we get $MP = 4$.

 We notice that $AMPC$ is a right trapezoid. Hence,
 \begin{align*} AC & = \sqrt{MP^2 - \left( AM - CP \right)^2} \\ & = \sqrt{15} . \end{align*}


 Therefore,
 \begin{align*} [XABY] & = \frac{1}{2} \left( AB + XY \right) AC \\ & = \frac{1}{2} \left( 12 + 24 \right) \sqrt{15} \\ & = 18 \sqrt{15}. \end{align*}


 Therefore, the answer is $18 + 15 = \boxed{\textbf{(033)}}$.

 



\textbf{Solution 2}

Notice that line $\overline{PQ}$ is the \href{https://artofproblemsolving.com/wiki/index.php?title=Radical\_axis}{radical axis} of circles $\omega_1$ and $\omega_2$. By the radical axis theorem, we know that the tangents of any point on line $\overline{PQ}$ to circles $\omega_1$ and $\omega_2$ are equal. Therefore, line $\overline{PQ}$ must pass through the midpoint of $\overline{AB}$, call this point M. In addition, we know that $AM=MB=6$ by circle properties and midpoint definition.

 Then, by Power of Point,

 
 \begin{align*} AM^2&=MP*MQ \\ 36&=MP(MP+5) \\ MP&=4 \\ \end{align*}


 Call the intersection point of line $\overline{A\omega_1}$ and $\overline{XY}$ be C, and the intersection point of line $\overline{B\omega_2}$ and $\overline{XY}$ be D. $ABCD$ is a rectangle with segment $MP=4$ drawn through it so that $AM=MB=6$, $CP=5$, and $PD=7$. Dropping the altitude from $M$ to $\overline{XY}$, we get that the height of trapezoid $XABY$ is $\sqrt{15}$. Therefore the area of trapezoid $XABY$ is

 
 \begin{align*} \frac{1}{2}\cdot(12+24)\cdot(\sqrt{15})=18\sqrt{15} \end{align*}


 Giving us an answer of $\boxed{033}$.

 



\textbf{Solution 3}

Refer to Solution 1.

 We let $AC=BD=x$ and the extension of $AC$ to the circle $\neq A$ as $E.$ By Power of a Point on point $C$ of circle $w_1$ we find 
 \[x\cdot{CE}=5\cdot5.\] 
 \[CE=\frac{25}{x}.\] We have diameter $AE = AC+CE=x+\frac{25}{x}.$ Therefore the radius of $w_1$ is $\frac{x+\frac{25}{x}}{2}=\frac{25+x^2}{2x} = O_1A = O_1P.$

 Similarly repeating this procedure on $w_2$ we find the radius of $w_2$ is $\frac{49+x^2}{2x} = O_2P = O_2B.$

 Next we solve for $O_1O_2$ in two ways. Let the perpendicular from $O_1$ to $BO_2$ intersect at $K$ we have $O_1K =AB = 12.$ We also have 
 \[O_2K = BO_2 - AO_1 =\frac{49+x^2}{2x}- \frac{25+x^2}{2x}=\frac{12}{x}.\]Therefore since $\triangle{O_1KO_2}$ is right, we have $(O_1O_2)^2 = (O_1K)^2+(O_2K)^2 = 12^2 + \frac{12}{x}^2 =144 + \frac{144}{x^2}.$

 For our second way, we let the midpoint of $PQ$ be $M.$ Note that $PM$ forms the right triangles $PO_1M$ and $PO_2M$ both of which share an leg of $PM$ or $\frac{5}{2}.$ Using Pythag we can solve for $O_1O_2.$

 \[O_1O_2 = O_1M+O_2M = \sqrt{(O_1P)^2 - (PM)^2}+\sqrt{(O_2P)^2 - (PM)^2}\]
 \[\sqrt{144 + \frac{144}{x^2}} = \sqrt{(\frac{25+x^2}{2x})^2 - (\frac{5}{2})^2}+\sqrt{(\frac{49+x^2}{2x})^2 - (\frac{5}{2})^2}\]
 \[\sqrt{144 + \frac{144}{x^2}} = \sqrt{\frac{(25+x^2)^2}{4x^2} - \frac{25}{4}}+\sqrt{\frac{(49+x^2)^2}{4x^2} - \frac{25}{4}}\]We let $x^2 = a$ to slightly simplify the equation,
 \[\sqrt{144 + \frac{144}{a}} = \sqrt{\frac{(25+a)^2}{4a} - \frac{25}{4}}+\sqrt{\frac{(49+a)^2}{4a} - \frac{25}{4}}\]
 \[12\sqrt{1+\frac{1}{a}} = \sqrt{\frac{a^2 + 25a + 625}{4a}} + \sqrt{\frac{a^2 +73a + 2401}{4a}}.\]
 \[24\sqrt{1+\frac{1}{a}} = \sqrt{\frac{a^2 + 25a + 625}{a}} + \sqrt{\frac{a^2 +73a + 2401}{a}}.\]
 \[24\sqrt{a+1} = \sqrt{a^2 + 25a + 625} + \sqrt{a^2 +73a + 2401}.\]
 \[\sqrt{a^2 + 25a + 625}=24\sqrt{a+1}-\sqrt{a^2 +73a + 2401}.\]
 \[a^2 + 25a + 625=576a+576+a^2+73a+2401-48\sqrt{(a+1)(a^2 +73a + 2401)}.\]
 \[48\sqrt{(a+1)(a^2 +73a + 2401)}=624a+2352.\]
 \[\sqrt{(a+1)(a^2 +73a + 2401)}=13a+49.\]
 \[a^3 + 74 a^2 + 2474 a + 2401=169 a^2 + 1274 a + 2401.\]
 \[a^3 -95a^2 + 1200a = 0\]
 \[a(a-15)(a-80)=0\]Thus the solutions are $a=0,15,80.$ Checking bounds $a=15$ is the only valid solution, which means $x=\sqrt{15}.$Finally to find the area of $XABY,$ we have the bases $XY=24$ and $AB=12$ and the height $x=\sqrt{15}$ therefore 
 \[[XABY]=\frac{1}{2}\cdot(12+24)\cdot\sqrt{15}=18\sqrt{15}.\]Giving us an answer of $18+15 = \boxed{033}.$



	\item (\textit{AIME II 2023}) Let $N$ be the number of ways to place the integers $1$ through $12$ in the $12$ cells of a $2 \times 6$ grid so that for any two cells sharing a side, the difference between the numbers in those cells is not divisible by $3.$ One way to do this is shown below. Find the number of positive integer divisors of $N.$

 \[\begin{array}{|c|c|c|c|c|c|} \hline \,1\, & \,3\, & \,5\, & \,7\, & \,9\, & 11 \\ \hline \,2\, & \,4\, & \,6\, & \,8\, & 10 & 12 \\ \hline \end{array}\]

 \textbf{Answer: }144



\textbf{Solution 1}

We replace the numbers which are $0 \pmod 3$ to $0$, $1 \pmod 3$ to $1$, and $2 \pmod 3$ to $2$.

 Then, the problem is equivalent to arranging four $0$'s, four $1$'s, and four $2$'s into the grid (and then multiplying by $4!^3$ to account for replacing the remainder numbers with actual numbers) such that no 2 of the same numbers are adjacent. 

 Then, the numbers which are vertically connected must be different. Two 1s must be connected with two 2s, and two 1s must be connected with two 0s (vertically), as if there are less 1s connected with either, then there will be 2s or 0s which must be connected within the same number.For instance if we did:

 
        \item  Three (1,2) pairs 
        \item  One (1,0) pair Then we would be left with one (0,2), and two remaining 0s which aren't supposed to be connected, so it is impossible. Similar logic works for why all 1s can't be connected with all 2s. Thus, we are left with the problem of re-arranging two (1,2) pairs, two (1,0) pairs, two (2,0) pairs. Notice that the pairs can be re-arranged horizontally in any configuration, but when two pairs are placed adjacent there is only one configuration for the rightmost pair to be set after the leftmost one has been placed.We have $\frac{6!}{2!2!2!}=90$ ways to horizontally re-arrange the pairs, with 2 ways to set the initial leftmost column. Thus, there are 180 ways to arrange the pairs. Accounting for the initial simplification of the problem with 1-12 $\rightarrow$ 0-3, we obtain the answer is: 
 \[2488320=2^{11}\cdot3^5\cdot5^1\]The number of divisors is: $12\cdot6\cdot2=\boxed{144}.$~SAHANWIJETUNGA

 



\textbf{Solution 2 (Archaeological)}

Let's carry out an archaeological study, that is, we will find the bones (the base) and "build up the meat."

 1. Let "bones" of number $X$ be $X \pmod 3.$ Then the “skeleton” of the original table is

 
 \[1 0 2 1 0 2\]
 \[2 1 0 2 1 0\]By condition, the table cannot have a column of two identical numbers (the difference of such numbers is a multiple of $3).$

 There are $4$ zeros, $4$ ones and $4$ twos in the table (the order of the numbers in the columns is not important).

 Therefore, there cannot be more than two identical columns (otherwise there will be four identical numbers in the remaining three, that is, the numbers in one column are the same).

 Any such table has exactly $2$ columns $(0,1), 2$ columns $(0,2)$ and $2$ columns $(1,2).$

 The number of possible tables of six elements of three types is 
 \[m = \frac {6!}{2! \cdot 2! \cdot 2!} = 2 \cdot 3^2 \cdot 5.\]The number of possible tables of six elements, if the order of the digits is important, is 
 \[M = 2^6 \cdot m = 2^7 \cdot 3^2 \cdot 5.\]2. We are looking for the total number of options.The column $(0,1)$ can contain the following $4^2 = 16$ pairs of numbers (the order is not important, it is already taken into account)
 \[(1,3), (1,6), (1,9), (1, 12), (4.3), (4.6), (4.9) (4.12), (7.3), (7.6), (7.9), (7.12), (11,3), (11,6), (11,9), (11,12).\]The second column $(0,1)$ can contain $3^2 = 9$ pairs of numbers (excluding the two that are in the first column).

 Similarly with the other two columns, i.e. the total number of choices 
 \[M \cdot 16 \cdot 9 \cdot 3 = 2^{11} \cdot 3^5 \cdot 5.\]Number of divisors 
 \[(11+1) \cdot (5+1) \cdot (1 + 1) = \boxed{144}.\]\textbf{vladimir.shelomovskii@gmail.com, vvsss}

 \textbf{Simpler Way to Finish} Note that there are $6!/(2!2!2!)$ ways for the possible column arrangements. Then, each of those columns have two numbers, which can be permuted either way, so we multiply by $2!$. Finally, we have $4$ numbers of each residue of $0, 1, 2 \pmod 3$. Thus, let us put those $4$ who are $0 \pmod 3$, and we get $4!$ ways to do this. Similarly, we can do the same for the 4 numbers each that are $1, 2 \pmod 3$, so we multiply by $(4!)^3$. Thus, the final answer is $\frac{6!}{2!2!2!}*2!*(4!)^3$ which leads us to the same solution as above.

 \textbf{mathboy282}

 LaTeX Edit: ~Bigbrain\_2009

 


	\item (\textit{AIME II 2023}) Find the number of collections of $16$ distinct subsets of $\{1,2,3,4,5\}$ with the property that for any two subsets $X$ and $Y$ in the collection, $X \cap Y \not= \emptyset.$

 \textbf{Answer: }081



\textbf{Solution 1}

Denote by $\mathcal C$ a collection of 16 distinct subsets of $\left\{ 1, 2, 3, 4, 5 \right\}$.Denote $N = \min \left\{ |S|: S \in \mathcal C \right\}$.

 Case 1: $N = 0$.

 This entails $\emptyset \in \mathcal C$.Hence, for any other set $A \in \mathcal C$, we have $\emptyset \cap A = \emptyset$. This is infeasible.

 Case 2: $N = 1$.

 Let $\{a_1\} \in \mathcal C$.To get $\{a_1\} \cap A \neq \emptyset$ for all $A \in \mathcal C$.We must have $a_1 \in \mathcal A$.

 The total number of subsets of $\left\{ 1, 2, 3, 4, 5 \right\}$ that contain $a_1$ is $2^4 = 16$.Because $\mathcal C$ contains 16 subsets.We must have $\mathcal C = \left\{ \{a_1\} \cup A : \forall \ A \subseteq \left\{ 1, 2, 3, 4, 5 \right\} \backslash \left\{a_1 \right\} \right\}$.Therefore, for any $X, Y \in \mathcal C$, we must have $X \cap Y \supseteq \{a_1\}$.So this is feasible.

 Now, we count the number of $\mathcal C$ in this case.We only need to determine $a_1$.Therefore, the number of solutions is 5.

 Case 3: $N = 2$.

 Case 3.1: There is exactly one subset in $\mathcal C$ that contains 2 elements.

 Denote this subset as $\left\{ a_1, a_2 \right\}$.We then put all subsets of $\left\{ 1, 2, 3, 4, 5 \right\}$ that contain at least three elements into $\mathcal C$, except $\left\{ a_3, a_4, a_5 \right\}$.This satisfies $X \cap Y \neq \emptyset$ for any $X, Y \in \mathcal C$.

 Now, we count the number of $\mathcal C$ in this case.We only need to determine $\left\{ a_1, a_2 \right\}$.Therefore, the number of solutions is $\binom{5}{2} = 10$.

 Case 3.2: There are exactly two subsets in $\mathcal C$ that contain 2 elements.

 They must take the form $\left\{ a_1, a_2 \right\}$ and $\left\{ a_1, a_3 \right\}$.

 We then put all subsets of $\left\{ 1, 2, 3, 4, 5 \right\}$ that contain at least three elements into $\mathcal C$, except $\left\{ a_3, a_4, a_5 \right\}$ and $\left\{ a_2, a_4, a_5 \right\}$.This satisfies $X \cap Y \neq \emptyset$ for any $X, Y \in \mathcal C$.

 Now, we count the number of $\mathcal C$ in this case.We only need to determine $\left\{ a_1, a_2 \right\}$ and $\left\{ a_1, a_3 \right\}$.Therefore, the number of solutions is $5 \cdot \binom{4}{2} = 30$.

 Case 3.3: There are exactly three subsets in $\mathcal C$ that contain 2 elements.They take the form $\left\{ a_1, a_2 \right\}$, $\left\{ a_1, a_3 \right\}$, $\left\{ a_1, a_4 \right\}$.

 We then put all subsets of $\left\{ 1, 2, 3, 4, 5 \right\}$ that contain at least three elements into $\mathcal C$, except $\left\{ a_3, a_4, a_5 \right\}$, $\left\{ a_2, a_4, a_5 \right\}$, $\left\{ a_2, a_3, a_5 \right\}$.This satisfies $X \cap Y \neq \emptyset$ for any $X, Y \in \mathcal C$.

 Now, we count the number of $\mathcal C$ in this case.We only need to determine $\left\{ a_1, a_2 \right\}$, $\left\{ a_1, a_3 \right\}$, $\left\{ a_1, a_4 \right\}$.Therefore, the number of solutions is $5 \cdot \binom{4}{3} = 20$.

 Case 3.4: There are exactly three subsets in $\mathcal C$ that contain 2 elements.They take the form $\left\{ a_1, a_2 \right\}$, $\left\{ a_1, a_3 \right\}$, $\left\{ a_2, a_3 \right\}$.

 We then put all subsets of $\left\{ 1, 2, 3, 4, 5 \right\}$ that contain at least three elements into $\mathcal C$, except $\left\{ a_3, a_4, a_5 \right\}$, $\left\{ a_2, a_4, a_5 \right\}$, $\left\{ a_1, a_4, a_5 \right\}$.This satisfies $X \cap Y \neq \emptyset$ for any $X, Y \in \mathcal C$.

 Now, we count the number of $\mathcal C$ in this case.We only need to determine $\left\{ a_1, a_2 \right\}$, $\left\{ a_1, a_3 \right\}$, $\left\{ a_2, a_3 \right\}$.Therefore, the number of solutions is $\binom{5}{3} = 10$.

 Case 3.5: There are exactly four subsets in $\mathcal C$ that contain 2 elements.They take the form $\left\{ a_1, a_2 \right\}$, $\left\{ a_1, a_3 \right\}$, $\left\{ a_1, a_4 \right\}$, $\left\{ a_1, a_5 \right\}$.

 We then put all subsets of $\left\{ 1, 2, 3, 4, 5 \right\}$ that contain at least three elements into $\mathcal C$, except $\left\{ a_3, a_4, a_5 \right\}$, $\left\{ a_2, a_4, a_5 \right\}$, $\left\{ a_1, a_4, a_5 \right\}$, $\left\{ a_2, a_3, a_4 \right\}$.This satisfies $X \cap Y \neq \emptyset$ for any $X, Y \in \mathcal C$.

 Now, we count the number of $\mathcal C$ in this case.We only need to determine $\left\{ a_1, a_2 \right\}$, $\left\{ a_1, a_3 \right\}$, $\left\{ a_1, a_4 \right\}$, $\left\{ a_1, a_5 \right\}$.Therefore, the number of solutions is 5.

 Putting all subcases together, the number of solutions is this case is $10 + 30 + 20 + 10 + 5 = 75$.

 Case 4: $N \geq 3$.

 The number of subsets of $\left\{ 1, 2, 3, 4, 5 \right\}$ that contain at least three elements is $\sum_{i=3}^5 \binom{5}{i} = 16$.Because $\mathcal C$ has 16 elements, we must select all such subsets into $\mathcal C$.Therefore, the number of solutions in this case is 1.

 Putting all cases together, the total number of $\mathcal C$ is $5 + 75 + 1 = \boxed{\textbf{(081) }}$.

 



\textbf{Solution 2}

Denote the $A$ as $\{ 1,2,3,4,5 \}$ and the collection of subsets as $S$. 

 
\textbf{Case 1: There are only sets of size $3$ or higher in $S$:  }Any two sets in $S$ must have at least one element common to both of them (since $3+3>5$). Since there are $16$ subsets of $A$ that have size $3$ or higher, there is only one possibility for this case.

 
\textbf{Case 2: There are only sets of size $2$ or higher in $S$:  }Firstly, there cannot be a no size $2$ set $S$, or it will be overcounting the first case.

 If there is only one such size $2$ set, there are $10$ ways to choose it. That size $2$ set, say $X$, cannot be in $S$ with $Y = A/X$ (a three element set). Thus, there are only $15$ possible size $3$ subsets that can be in $S$, giving us $10$ for this case.

 If there are two sets with size $2$ in $S$, we can choose the common elements of these two subsets in $5$ ways, giving us a total of $5\cdot 6 = 30$.

 If there are three sets with size $2$, they can either share one common element, which can be done in $5 \cdot 4 = 20$ ways, or they can share pairwise common elements (sort of like a cycle), which can be done $\binom{5}{2} = 10$ ways. In total, we have $30$ possibilities.

 If there are four sets with size $2$, they all have to share one common element, which can be done in $5\cdot 1$ ways.

 Thus, summing everything up, this will give us $75$ possible sets $S$

 
\textbf{Case 3: There is a set with size $1$ in $S$:  }Notice that be at most one size $1$ subset. There are $5$ ways to choose that single element set. Say it's $\{ 1\}$. All other subsets in $S$ must have a $1$ in them, but there are only $2^4-1 = 15$ of them. Thus this case yields $5\cdot 1 = 5$ possibilities.

 
Thus, the total number of sets $S$ would be $1+75+5 = \boxed{\textbf{(081)}}$.

 



\textbf{Solution 3}

Firstly, there cannot be two subsets with cardinality 1, or they will not intersect.

 If there is one subset $A$ with cardinality $1$; let the element in $A$ be $a$, then there are $2^4=16$ subsets that do not include $a$ so they do not work. Every remaining subsets $S$ will have $a$ as an element so $S\cap{A}\geq1$, since we just excluded all that do not. Since there are 15 subsets left, all are forced into the group of 16 subsets, so we just choose the number of $a$ to determine the set $A$. $A\in\{1,2,3,4,5\}$ so there are 5 ways.

 For the rest of the cases, we assume there are no sets $A$ with cardinality 1.  Notice that the only way to "violate" the condition is to have subsets $X$ and $Y$ with cardinalities 2 and 3 in some order. Otherwise, by the Pigeonhole Principle, if two sets both have cardinalities more than 3, they are bound to have one element of intersection. Say a set $S$ has $|S|=2$, then there is clearly only one set $s$ that will make $S\cap{s}=0$. By our previous claim, all other subsets that have cardinality $c\geq{3}$ will work.

 Now if we generalize a bit: If a subset has $N$ 2-element subsets which belong to set $M$, then there are exactly $N$ subsets with cardinality 3 that don't work. Therefore, the number of "violating subsets" are all subsets with cardinality $c\leq{1}$, all 2-element subsets that are not in $M$, and all corresponding cardinality 3 subsets. Subtracting from the total 32 subsets, we get that $32-(1+5+(10-N)+N)=16$ subsets that do work. This includes all subsets in $M$, so the remaining non-violating subsets are forced. This is equivalent now to choosing $N$ 2 element subsets.

 Following casework on the number of 2-element subsets:

 If $N=1$: There are $\binom{5}{2}=10$ ways.

 If $N=2$: There are $5$ ways to choose the intersection between the 2 sets (remember they have to have at least one element of intersection) and $\binom{4}{2}=6$ ways to choose the distinct elements in the subsets, so there are $5*6=30$ ways.

 If $N=3$: It can be a cycle, where WLOG let the elements be $a,b,c$ so the sets are $\{a,b\}$, $\{b,c\}$, and $\{c,a\}$. This is just $\binom{5}{3}=10$. Alternatively, it can also be the case where all sets share one element. There are 5 ways to choose this element and $\binom{4}{2}=6$ ways to choose the remaining elements to assign to each set. There are $30$ ways.

 If $N=4$: By the Pigeonhole Principle, the only way all pairwise sets have at one common intersection is if all share one element in common. There are 5 ways to choose this element and the remaining numbers are forced. There are 5 ways.

 $N\geq{5}$ does not provide any valid cases since to have all pairwise elements to intersect one element, they must be the same element by the Pigeonhole Principle, but there are not enough subsets.

 If $N=0$, then there is only one way since $\binom{5}{3}+\binom{5}{4}+\binom{5}{5}=16$.

 Adding all cases yields $5+10+30+30+5+1=\boxed{81}$ ways!

 -Magnetoninja

 


	\item (\textit{AIME II 2023}) In $\triangle ABC$ with side lengths $AB = 13,$ $BC = 14,$ and $CA = 15,$ let $M$ be the midpoint of $\overline{BC}.$ Let $P$ be the point on the circumcircle of $\triangle ABC$ such that $M$ is on $\overline{AP}.$ There exists a unique point $Q$ on segment $\overline{AM}$ such that $\angle PBQ = \angle PCQ.$ Then $AQ$ can be written as $\frac{m}{\sqrt{n}},$ where $m$ and $n$ are relatively prime positive integers. Find $m+n.$

 \textbf{Answer: }247



\textbf{Solution 1}

$\textbf{Claim: BPCQ is a parallelogram}$
$\textit{Proof}$: Let $Q'$ be such that $BPCQ'$ is a parallelogram. Then, $\angle PBQ' = \angle PCQ'$ and $Q' \in \overline{PM}$, since diagonals bisect. Thus, $Q'=Q$.

 Now, by Stewart's or Apollonius', $AM = \sqrt{148}$. Then, by power of a point, 
 \[AM \cdot MP = BM \cdot CM = 49 \implies PM = \frac{49}{AM}\] Thus, noting $PM = QM$, 
 \[AQ = AM - MQ = \frac{AM^2-49}{AM} = \frac{99}{\sqrt{148}} \implies \boxed{247}\]

 -sillybone

 



\textbf{Solution 2}

Because $M$ is the midpoint of $BC$, following from the Stewart's theorem, $AM = 2 \sqrt{37}$.

 Because $A$, $B$, $C$, and $P$ are concyclic, $\angle BPA = \angle C$, $\angle CPA = \angle B$.

 Denote $\theta = \angle PBQ$.

 In $\triangle BPQ$, following from the law of sines,
 \[ \frac{BQ}{\sin \angle BPA} = \frac{PQ}{\sin \angle PBQ} \]

 Thus,
 \[ \frac{BQ}{\sin C} = \frac{PQ}{\sin \theta} . \hspace{1cm} (1) \]

 In $\triangle CPQ$, following from the law of sines,
 \[ \frac{CQ}{\sin \angle CPA} = \frac{PQ}{\sin \angle PCQ} \]

 Thus,
 \[ \frac{CQ}{\sin B} = \frac{PQ}{\sin \theta} . \hspace{1cm} (2) \]

 Taking $\frac{(1)}{(2)}$, we get
 \[ \frac{BQ}{\sin C} = \frac{CQ}{\sin B} \]

 In $\triangle ABC$, following from the law of sines,
 \[ \frac{AB}{\sin C} = \frac{AC}{\sin B} . \hspace{1cm} (3) \]

 Thus, Equations (2) and (3) imply
 \begin{align*} \frac{BQ}{CQ} & = \frac{AB}{AC} \\ & = \frac{13}{15} . \hspace{1cm} (4) \end{align*}


 


 Next, we compute $BQ$ and $CQ$.

 We have
 \begin{align*} BQ^2 & = AB^2 + AQ^2 - 2 AB\cdot AQ \cos \angle BAQ \\ & = AB^2 + AQ^2 - 2 AB\cdot AQ \cos \angle BAM \\ & = AB^2 + AQ^2 - 2 AB\cdot AQ \cdot \frac{AB^2 + AM^2 - BM^2}{2 AB \cdot AM} \\ & = AB^2 + AQ^2 -  AQ \cdot \frac{AB^2 + AM^2 - BM^2}{AM} \\ & = 169 + AQ^2 - \frac{268}{2 \sqrt{37}} AQ .  \hspace{1cm} (5) \end{align*}


 We have
 \begin{align*} CQ^2 & = AC^2 + AQ^2 - 2 AC\cdot AQ \cos \angle CAQ \\ & = AC^2 + AQ^2 - 2 AC\cdot AQ \cos \angle CAM \\ & = AC^2 + AQ^2 - 2 AC\cdot AQ \cdot \frac{AC^2 + AM^2 - CM^2}{2 AC \cdot AM} \\ & = AC^2 + AQ^2 -  AQ \cdot \frac{AC^2 + AM^2 - CM^2}{AM} \\ & = 225 + AQ^2 - \frac{324}{2 \sqrt{37}} AQ .  \hspace{1cm} (6) \end{align*}


 Taking (5) and (6) into (4), we get $AQ = \frac{99}{\sqrt{148}}$

 Therefore, the answer is $99 + 148 = \boxed{\textbf{(247) }}$




\textbf{Solution 3}

Define $L_1$ to be the foot of the altitude from $A$ to $BC$. Furthermore, define $L_2$ to be the foot of the altitude from $Q$ to $BC$. From here, one can find $AL_1=12$, either using the 13-14-15 triangle or by calculating the area of $ABC$ two ways. Then, we find $BL_1=5$ and $L_1C = 9$ using Pythagorean theorem. Let $QL_2=x$. By AA similarity, $\triangle{AL_1M}$ and $\triangle{QL_2M}$ are similar. By similarity ratios, 
 \[\frac{AL_1}{L_1M}=\frac{QL_2}{L_2M}\]
 \[\frac{12}{2}=\frac{x}{L_2M}\]
 \[L_2M = \frac{x}{6}\]Thus, $BL_2=BM-L_2M=7-\frac{x}{6}$. Similarly, $CL_2=7+\frac{x}{6}$. Now, we angle chase from our requirement to obtain new information.
 \[\angle{PBQ}=\angle{PCQ}\]
 \[\angle{QCM}+\angle{PCM}=\angle{QBM}+\angle{PBM}\]
 \[\angle{QCL_2}+\angle{PCM}=\angle{QBL_2}+\angle{PBM}\]
 \[\angle{PCM}-\angle{PBM}=\angle{QBL_2}-\angle{QCL_2}\]
 \[\angle{MAB}-\angle{MAC}=\angle{QBL_2}-\angle{QCL_2}\text{(By inscribed angle theorem)}\] 
 \[(\angle{MAL_1}+\angle{L_1AB})-(\angle{CAL_1}-\angle{MAL_1})=\angle{QBL_2}-\angle{QCL_2}\]
 \[2\angle{MAL_1}+\angle{L_1AB}-\angle{CAL_1}=\angle{QBL_2}-\angle{QCL_2}\]Take the tangent of both sides to obtain 
 \[\tan(2\angle{MAL_1}+\angle{L_1AB}-\angle{CAL_1})=\tan(\angle{QBL_2}-\angle{QCL_2})\]By the definition of the tangent function on right triangles, we have $\tan{MAL_1}=\frac{7-5}{12}=\frac{1}{6}$, $\tan{CAL_1}=\frac{9}{12}=\frac{3}{4}$, and $\tan{L_1AB}=\frac{5}{12}$. By abusing the tangent angle addition formula, we can find that 
 \[\tan(2\angle{MAL_1}+\angle{L_1AB}-\angle{CAL_1})=\frac{196}{2397}\]By substituting $\tan{\angle{QBL_2}}=\frac{6x}{42-x}$, $\tan{\angle{QCL_2}}=\frac{6x}{42+x}$ and using tangent angle subtraction formula we find that 
 \[x=\frac{147}{37}\]Finally, using similarity formulas, we can find 
 \[\frac{AQ}{AM}=\frac{12-x}{x}\]. Plugging in $x=\frac{147}{37}$ and $AM=\sqrt{148}$, we find that 
 \[AQ=\frac{99}{\sqrt{148}}\] Thus, our final answer is $99+148=\boxed{247}$. ~sigma

 



\textbf{Solution 4}

It is clear that $BQCP$ is a parallelogram. By Stewart's Theorem, $AM=\sqrt{148}$, POP on $M$ tells $PM=\frac{49}{\sqrt{148}}$

 As $QM=PM, AQ=AM-PM=\frac{99}{\sqrt{148}}$ leads to $\boxed{247}$
$\textbf{NOTE: Why BQCP is a parallelogram}$

 It's not actually immediately clear why this is the case. There are two ways to easily show this:

 $\textbf{Competition solution:}$

 Notice that the problem statement tells us that point Q is $\textit{unique.}$ EVERY piece of information in the problem statement is intentional, so we should try to use this to our benefit. None of the other solutions do, which is why they are more complicated than they need be.

 Consider point Q' s.t. $Q'M = MP$. Obviously, $\angle Q'CP$ and $\angle Q'BP$ are equal - we have perfect symmetry along line $AP$. Moreover, $BQ'CP$ is a parallelogram as its diagonals bisect each other. Since point $Q$ is unique, we know that $Q' \textit{is } Q$. Thus $BQCP$ is a parallelogram. $\blacksquare$$\newline$$\textbf{Rigorous proof (not recommended for competition scenario):}$Consider any quadrilateral $ABCD$ whose diagonals intersect at $O$ s.t. $AO = OC$ and $\angle BAD = \angle BCD$. We will prove that $ABCD$ is $\textit{either a \textbf{parallelogram} or a \textbf{kite}}$. 

 (Note that in our problem, since $AP$ and $BC$ are not orthogonal, ($ABC$ isn't isosceles) this is enough to show that $BQCP$ is a parallelogram).$\newline$-- By same base/same altitude, $[ABO] = [CBO]$ and $[ADO] = [CDO] \implies [ABD] = [ABO] + [ADO] = [CBO] + [CDO] = [CBD]. \newline$

 Therefore: $\frac{1}{2} sin(\angle BAD) \overline{AB} \times \overline{AD} = \frac{1}{2} sin(\angle BCD) \overline{CB} \times \overline{CD}.$ Since $\angle BAD = \angle BCD$, this reduces to $\overline{AB} \times \overline{AD} = \overline{CB} \times \overline{CD}. (E.1) \newline$

 Let $AB = x$ and $AD = y$. Then, by $(E.1)$, $CB = kx$ and $CD = \frac{y}{k}$ for some real $k$. Then by LoC on $\triangle BAD$ and $\triangle BCD$:$x^{2} + y^{2} - 2xy cos(\angle BAD) = \overline{BD} = x^{2}k^{2} + \frac{y^{2}}{k^{2}} - 2xy cos(\angle BCD) \newline \implies x^{2} + y^{2} = x^{2}k^{2} + \frac{y^{2}}{k^{2}} \newline \implies (y^{2} - x^{2}k^{2})(k^{2} - 1) = 0.\newline$

 -- $y^{2} - x^{2}k^{2} = 0 \implies y = kx \implies AD = BC$ and $AB = CD \implies$ $ABCD$ is a parallelogram.

 -- $k^{2} - 1 = 0 \implies k = 1$ ($k$ cannot be $-1$; no negative sided polygons here!) $\implies AB = CB$ and $AD = CD \implies$ $ABCD$ is a kite. $\square$. ~Mathavi

 



\textbf{Solution 5 (LOS+ coordbash)}

First, note that by Law of Sines, $\frac{\sin(\angle{QBP})}{QP}=\frac{\sin(\angle{BPQ})}{BQ}$ and that $\frac{\sin(\angle{QCP})}{QP}=\frac{\sin(\angle{QPC})}{QP}$. Equating the 2 expressions, you get that $\frac{\sin(\angle{BPQ})}{BQ}=\frac{\sin(\angle{QPC})}{QP}$. Now drop the altitude from $A$ to $BC$. As it is commonly known that the dropped altitude forms a $5-12-13$ and a $9-12-15$ triangle, you get the measures of $\angle{ABC}$ and $\angle{ACB}$ respectively, which are $\arcsin(\frac{12}{13})$ and $\arcsin(\frac{4}{5})$. However, by the inscribed angle theorem, you get that $\angle{BPQ}=\arcsin(\frac{4}{5})$ and that $\angle{QPC}=\arcsin(\frac{12}{13})$, respectively. Therefore, by Law of Sines (as previously stated) $\frac{BQ}{CQ}=\frac{13}{15}$.

 Now commence coordbashing. Let $B$ be the origin, and $A$ be the point $(5,12)$. As $AP$ passes through $A$, which is $(5,12)$, and $M$, which is $(7,0)$, it has the equation $-6x+42$, so therefore a point on this line can be written as $(x,42-6x)$. As we have the ratio of the lengths, which prompts us to write the lengths in terms of the distance formula, we can just plug and chug it in to get the ratio $\frac{\sqrt{37x^2-504x+1764}}{\sqrt{37x^2-532x+1960}}=\frac{13}{15}$. This can be squared to get $\frac{37x^2-504x+1764}{37x^2-532x+1960}=\frac{169}{225}$. This can be solved to get a solution of $x=\frac{469}{74}$, and an extraneous solution of $5$ which obviously doesn’t work. 

 Plugging $x$ into the line equation gets you $y=\frac{147}{37}$. The distance between this point and $A$, which is $(5,12)$ is $\sqrt{\frac{9801}{148}}$, or simplified to $\frac{99}{\sqrt{148}}\Longrightarrow99+148=\boxed{247}$




\textbf{Solution 6 (similar to 3)}


                    \begin{figure}[H]
                        \centering
                        \includegraphics[width=0.5\textwidth]{Solution/AIME_II_2023_12_Solution_Figure_1.png}
                    \end{figure}
                    We use the law of Cosine and get 
 \[AB^2 = AM^2 + BM^2 - 2 AM \cdot BM \cos \angle AMB,\] 
 \[AC^2 = AM^2 + CM^2 + 2 AM \cdot CM \cos \angle AMB \implies\] 
 \[AM^2 = \frac {AB^2 + AC^2}{2}- BM^2 = \sqrt{148} \approx 12.\]We use the power of point $M$ with respect circumcircle $\triangle ABC$ and get 
 \[AM \cdot MP = BM \cdot CM = BM^2 \implies\]
 \[PM = \frac {49}{\sqrt {148}} \approx \frac {48}{12} \approx 4 < AM.\]It is clear that if $Q = P,$ then $\angle PBQ = \angle PCQ = 0 \implies$ 

 if $Q$ is symmetric to $P$ with respect $M$ then $\angle PBQ = \angle PCQ.$

 There exists a <i>\textbf{unique}</i> point $Q$ on segment $\overline{AM}, PM < AM \implies$
 \[PQ = AM - PM = \frac{99}{\sqrt{148}} \implies \boxed{247}.\]\textbf{vladimir.shelomovskii@gmail.com, vvsss}

 



\textbf{Solution 7}

We will first find $AM$ by Stewart's Theorem. Hence:

 $15^{2} \cdot 7 + 13^{2} \cdot 7 = 14 \cdot 7 \cdot 7 + d^{2} \cdot 14$

 Where $d = AM$. Now simplifying:

 $225 + 169 = 98 + 2d^{2}$

 Which means:

 $d^{2} = 148$ or $d = \sqrt{148}$.

 Now we need to find $PM = QM$ because $M$ is the midpoint of both $BC$ and $PQ$. In fact, we will find $PM$ and the answer will just be $AM - PM = \sqrt{148} - PM$. 

 Let $BP = x$ and $PC = y$ and $PM = z$. Note that we can again apply Stewart's Theorem on triangle $BPC$
$7x^{2} + 7y^{2} = 14 \cdot 7 \cdot 7 + 14z^{2}$

 This yields:

 $x^{2} + y^{2} = 98 + 2z^{2}$

 And now, it remains to find $x$ and $y$ which can easily find using similar triangles and arc lengths. We will show triangle $BPM$ is similar to triangle $AMC$ and triangle $PMC$ is similar to triangle $ABM$. Now notice $\angle MBP = \angle MAC$ because they both intercept the same circular arc $PC$. Next, $\angle BMP = \angle AMC$ by vertical angles. Hence, we have an Angle-Angle similarity and thus triangles $BPM$ and $AMC$ are similar. We could do the same for triangle $PMC$. Note that $\angle MCP = \angle BAM$ because they both intercept the same circular arc $BP$ and $\angle CMP = \angle AMB$ by vertical angles. Hence, we again have Angle-Angle similarity and thus triangles $PMC$ is similar to triangle $ABM$. Now, we bring out the similarity ratios to solve for $x$ and $y$ and finish the problem. Note that we can solve for $x$ and $y$ by noticing $\frac{BP}{AC} = \frac{BM}{AM}$ and $\frac{PC}{AB} = \frac{MC}{AM}$. Hence, we can substitute values to get

 $\frac{x}{15} = \frac{7}{\sqrt{148}}$ and $\frac{y}{13} = \frac{7}{\sqrt{148}}$. Solving, $x = \frac{105}{\sqrt{148}}$ and $y = \frac{91}{\sqrt{148}}$. Recall that we found $x^{2} + y^{2} = 98 + 2z^{2}$. Now, we can solve:

 $\frac{105^{2}}{148} + \frac{91^{2}}{148} = 98 + 2z^{2}$

 To simplify this, note that we have to solve:

 $\frac{105^{2} + 91^{2}}{148} = 98 + 2z^{2}$

 Let's focus on $105^{2} + 91^{2}$ for now. Interestingly enough, both numbers are centered around $98$ which means notice $105 = 98 + 7$ and $91 = 98 - 7$. We can simplify this expression by using this fact to get $(98 + 7)^{2} + (98 - 7)^{2} = 98^{2} + 49 + 98^{2} + 49 = 2 \cdot 98^{2} + 98$.

 Now, we have:

 $\frac{2 \cdot 98^{2} + 98}{148} = 98 + 2z^{2}$

 Note that now we have:

 $\frac{ 2 \cdot 98 \cdot 98 + 98 - 98 \cdot 148}{148} = 2z^{2}$

 This yields:

 $\frac{98(196 + 1 - 148)}{148} = 2z^{2}$

 Now we have:

 $\frac{98 \cdot 49}{148} = 2z^{2}$

 We can cancel out the 2's and simplify to be left with:

 $\frac{49}{\sqrt{148}} = z = PM$.

 Finally, the answer is $AM - PM = \sqrt{148} - \frac{49}{\sqrt{148}} = \frac{148 - 49}{\sqrt{148}} = \frac{99}{\sqrt{148}}$ which yields a final answer to $99 + 148 = \boxed{247}$.

 $\textbf{Supplement: Proof of why M is also the midpoint of PQ}$

 In this solution, I claimed that $M$ is also the midpoint of $PQ$. In fact, we could prove this by noting that because $\angle CQM$ and and $\angle MBP$ share the same circular arc $PC$, these angles are congruent. Note that $\angle BMP = \angle QMC$ by vertical angles and note that $BM = CM = 7$. We have Angle-Side-Angle congruence that hence triangle $BMP$ is congruent to triangle $QMC$ and hence $QM = MP$ proving the claim. You could also refer to Solution 3's supplements of the proofs that $BQPC$ is actually a parallelogram.

 


	\item (\textit{AIME II 2023}) Let $A$ be an acute angle such that $\tan A = 2 \cos A.$ Find the number of positive integers $n$ less than or equal to $1000$ such that $\sec^n A + \tan^n A$ is a positive integer whose units digit is $9.$

 \textbf{Answer: }167



\textbf{Solution 1}

Denote $a_n = \sec^n A + \tan^n A$.For any $k$, we have
 \begin{align*} a_n & = \sec^n A + \tan^n A \\ & = \left( \sec^{n-k} A + \tan^{n-k} A \right) \left( \sec^k A + \tan^k A \right) - \sec^{n-k} A \tan^k A - \tan^{n-k} A \sec^k A \\ & = a_{n-k} a_k - 2^k \sec^{n-k} A \cos^k A - 2^k \tan^{n-k} A \cot^k A \\ & = a_{n-k} a_k - 2^k a_{n-2k} . \end{align*}


 Next, we compute the first several terms of $a_n$.

 By solving equation $\tan A = 2 \cos A$, we get $\cos A = \frac{\sqrt{2 \sqrt{17} - 2}}{4}$.Thus, $a_0 = 2$, $a_1 = \sqrt{\sqrt{17} + 4}$, $a_2 = \sqrt{17}$, $a_3 = \sqrt{\sqrt{17} + 4} \left( \sqrt{17} - 2 \right)$, $a_4 = 9$.

 In the rest of analysis, we set $k = 4$.Thus,
 \begin{align*} a_n & = a_{n-4} a_4 - 2^4 a_{n-8}  \\ & = 9 a_{n-4} - 16 a_{n-8} . \end{align*}


 Thus, to get $a_n$ an integer, we have $4 | n$.In the rest of analysis, we only consider such $n$. Denote $n = 4 m$ and $b_m = a_{4n}$.Thus,
 \begin{align*} b_m & = 9 b_{m-1} - 16 b_{m-2} \end{align*}
with initial conditions $b_0 = 2$, $b_1 = 9$.

 To get the units digit of $b_m$ to be 9, we have
 \begin{align*} b_m \equiv -1 & \pmod{2} \\ b_m \equiv -1 & \pmod{5} \end{align*}


 Modulo 2, for $m \geq 2$, we have
 \begin{align*} b_m & \equiv 9 b_{m-1} - 16 b_{m-2} \\ & \equiv b_{m-1} . \end{align*}


 Because $b_1 \equiv -1 \pmod{2}$, we always have $b_m \equiv -1 \pmod{2}$ for all $m \geq 2$.

 Modulo 5, for $m \geq 5$, we have
 \begin{align*} b_m & \equiv 9 b_{m-1} - 16 b_{m-2} \\ & \equiv - b_{m-1} - b_{m-2} . \end{align*}


 We have $b_0 \equiv 2 \pmod{5}$, $b_1 \equiv -1 \pmod{5}$,$b_2 \equiv -1 \pmod{5}$, $b_3 \equiv 2 \pmod{5}$, $b_4 \equiv -1 \pmod{5}$, $b_5 \equiv -1 \pmod{5}$, $b_6 \equiv 2 \pmod{5}$.Therefore, the congruent values modulo 5 is cyclic with period 3.To get $b_m \equiv -1 \pmod{5}$, we have $3 \nmid m \pmod{3}$.

 From the above analysis with modulus 2 and modulus 5, we require $3 \nmid m \pmod{3}$.

 For $n \leq 1000$, because $n = 4m$, we only need to count feasible $m$ with $m \leq 250$.The number of feasible $m$ is
 \begin{align*} 250 - \left\lfloor \frac{250}{3} \right\rfloor & = 250 - 83 \\ & = \boxed{\textbf{(167) }} . \end{align*}





\textbf{Solution 2 (Simple)}


 \[\tan A = 2 \cos A \implies \sin A = 2 \cos^2 A \implies  \sin^2 A + \cos^2 A = 4  \cos^4 A + \cos^2 A = 1\]
 \[\implies   \cos^2 A = \frac {\sqrt {17} - 1}{8}.\]
 \[c_n = \sec^n A + \tan^n A = \frac {1}{\cos^n A} + 2^n \cos^n A = (4\cos^2 A +1)^{\frac {n}{2}}+(4 \cos^2 A)^{\frac {n}{2}} =\]
 \[= \left(\frac {\sqrt {17} + 1}{2}\right)^{\frac {n}{2}}+ \left(\frac {\sqrt {17} - 1}{2}\right)^{\frac {n}{2}}.\]

 It is clear, that $c_n$ is not integer if $n \ne 4k, k > 0.$

 Denote $x = \frac {\sqrt {17} + 1}{2}, y = \frac {\sqrt {17} - 1}{2} \implies$
 \[x \cdot y = 4, x + y = \sqrt{17}, x - y = 1 \implies x^2 + y^2 = (x - y)^2 + 2xy = 9 = c_4.\]

 \[c_8 = x^4 + y^4 = (x^2 + y^2)^2 - 2x^2 y^2 = 9^2 - 2 \cdot 16 = 49.\]
 \[c_{4k+4} = x^{2k+2} + y^{2k+2} = (x^{2k} + y^{2k})(x^2+y^2)- (x^2 y^2)(x^{2k-2}+y^{2k-2}) = 9 c_{4k}- 16 c_{4k – 4} \implies\]
 \[c_{12} = 9 c_8 - 16 c_4 = 9 \cdot 49 - 16 \cdot 9 = 9 \cdot 33 = 297.\]
 \[c_{16} = 9 c_{12} - 16 c_8 = 9 \cdot 297 - 16 \cdot 49 = 1889.\]
 \[c_{12m + 4} \pmod{10} = 9 \cdot c_{12m} \pmod{10} - 16 \pmod{10} \cdot c_{12m - 4}  \pmod{10} =\]
 \[= (9 \cdot 7 - 6 \cdot 9)  \pmod{10} = (3 - 4)  \pmod{10} = 9.\]
 \[c_{12m + 8}\pmod{10} = 9 \cdot c_{12m+4} \pmod{10} - 16 \pmod{10} \cdot c_{12m } \pmod{10} =\]
 \[= (9 \cdot 9 - 6 \cdot 7) \pmod{10} = (1 - 2)\pmod{10} = 9.\]
 \[c_{12m + 12} \pmod{10} = 9 \cdot c_{12m + 8} \pmod{10} - 16 \pmod{10} \cdot c_{12m + 4} \pmod{10} =\]
 \[= (9 \cdot 9 - 6 \cdot 9) \pmod{10} = (1 - 4) \pmod{10} = 7 \implies\]

 The condition is satisfied iff $n = 12 k + 4$ or $n = 12k + 8.$

 If $n \le N$ then the number of possible n is $\left\lfloor \frac{N}{4} \right\rfloor - \left\lfloor \frac{N}{12} \right\rfloor.$

 For $N = 1000$ we get $\left\lfloor \frac{1000}{4} \right\rfloor - \left\lfloor \frac{1000}{12} \right\rfloor = 250 - 83 = \boxed{167}.$

 \textbf{vladimir.shelomovskii@gmail.com, vvsss}

 



\textbf{Solution 3 (Characteristic equations)}

Using the method described in Solution 2, we can easily find $\cos A = \sqrt{\frac{\sqrt {17} - 1}{8}}$$\implies \tan A = 2\cos A = \sqrt{\frac{\sqrt{17} - 1}{2}}$

 Also, $\sec A = \frac{1}{\cos A} = \sqrt{\frac{8}{\sqrt {17} - 1}} = \sqrt{\frac{\sqrt{17} + 1}{2}}$

 Let $a = \sec^2A$, $b = \tan^2A \implies \sec^nA + \tan^nA = a^{n/2} + b^{n/2}$ 

 Notice $a$, $b$ are the roots of $x^2 - \sqrt {17}x + 4 = 0 \implies$$x^2 = \sqrt {17}x - 4$

 Define a sequence, $c_i = a^i + b^i \implies$ the recurrence formula for the sequence $c_i$ is $c_{n+1} = \sqrt{17}c_n - 4c_{n-1}$

 Computing all the even indexed terms $\pmod{10}$:\begin{align*}c_1 &= \sqrt{17} \pmod{10} \\c_2 &= 9 \pmod{10} \\c_3 &= 5\sqrt{17} \pmod{10} \\c_4 &= 49 \equiv 9 \pmod{10} \\c_{5} &= 29\sqrt{17} \equiv 9\sqrt{17} \pmod{10} \\c_{6} &= 297 \equiv 7 \pmod{10} \\c_{7} &= 181\sqrt{17} \equiv \sqrt{17} \pmod{10} \\c_{8} &= 1889 \equiv 9 \pmod{10} \\\end{align*}
It is clear this sequence repeats itself every 6 terms with 2 positive integers with a units digit of 9 every cycle.

 It also can be easily shown that any odd n wouldn't work because the odd n's would just be another sequence starting with $\sqrt{\sqrt{17} + 4}$ which will never be a positive integer.

 For $n \le 1000$, $\frac{n}{2} \le 500$ giving a total of $\left\lfloor \frac{1000}{6} \right\rfloor + 1 = 166 + 1 = \boxed{167}$

 \textbf{Note} A key idea in this solution is realizing that to calculate values of $a^n+b^n$ is difficult directly, so we try to think about it recursively.

 We then see how we can go from $a^n+b^n$ to $a^{n+1}+b^{n+1}$, and learn that $(a^n+b^n)(a+b)-a^nb-ab^n=a^{n+1}b^{n+1}$. 

 So, letting $a^i+b^i=s_i$, $s_i=(a+b)s_{i-1}-(ab)s_{i-2}$ where we set $n+1$ as $i$, and $(a^n+b^n)=s_{i-1}$, as well as $(a^nb+ab^n)=(ab)(a^{n-1}+b^{n-1})=(ab)s_{i-2}$.

 



\textbf{Solution 4(Rigorous)}

Since $A$ is acute, everything is positive so we don't need to worry about plus or minus. Let $x = \tan(A)$ and note that since $\tan^2(A) + 1 = \sec^2(A) \implies \frac{1}{x^2 + 1} = \cos^2(A)$ thus $\tan(A) = 2 \cos(A) \implies \tan^2(A) = 4\cos^2(A)$ so

 $x^2 = \frac{4}{x^2 + 1} \implies x^4 + x^2 - 4 = 0 \implies x^2 = \frac{-1 \pm \sqrt{1 + 16}}{2}$ but we can disregard the minus case since that would imply $x^2 < 0$ implying imaginary solutions, clearly absurd. So $x^2 = \frac{\sqrt{17} - 1}{2}$ and thus $\sec(A) = \sqrt{x^2 + 1} = \sqrt{\frac{\sqrt{17} + 1}{2}}$ and thus $\sec^n(A) + \tan^n(A)$ turns into

 $(\sqrt{\frac{\sqrt{17} + 1}{2}})^n + (\sqrt{\frac{\sqrt{17} - 1}{2}})^n$.

 Claim The above expression is only an integer if $n$ is a multiple of $4$

 Proof Let $a = \frac{\sqrt{17} + 1}{2}, b = \frac{\sqrt{17} - 1}{2}$ so the above expression is just $a^{\frac{n}{2}} + b^{\frac{n}{2}}$. We have $a + b = \sqrt{17}, ab = \frac{17 - 1}{4} = 4$ thus $a, b$ are the roots of the polynomial $x^2 - x\sqrt{17} + 4 = 0$ by Vieta's. Also, note that for $a^{\frac{n}{2}} + b^{\frac{n}{2}}$ we must have $n$ to be even otherwise we would have thins like $\sqrt{a}$ and $\sqrt{b}$ involved and you can check there's no way to simplify that so if $n$ were to be odd, being an integer would be impossible. Thus, let $n = 2k$. We will let $S_k = a^k + b^k$. Then, by Characteristic Polynomial, we get from the quadratic that

 $S_k = \sqrt{17} S_{k - 1} - 4S_{k - 2}$ and we manually know that $S_0 = a^0 + b^0 = 2, S_1 = a + b = \sqrt{17}$ so then we compute

 $S_2 = \sqrt{17} S_1 - 4 S_0 = \sqrt{17} \cdot \sqrt{17} - 4 \cdot 2 = 17 - 8 = 9$
$S_3 = \sqrt{17} S_2 - 4 S_1 = \sqrt{17} \cdot 9 - 4 \cdot \sqrt{17} = 5\sqrt{17}$
$S_4 = \sqrt{17} S_3 - 4 S_2 = \sqrt{17} \cdot 5 \sqrt{17} - 4 \cdot 9 = 5 \cdot 17 - 36$ obviously an integer.

 We note that all even terms gives an integer so we suspect $k$ also has to be an even integer. We need to prove this.We will use induction.

 Claim $S_{2m}$ is an integer, $S_{2m + 1}$ is $c_m \cdot \sqrt{17}$ for some integer $c_m$

 Proof The base case is $m = 0$ which implies $S_{2m} = S_0 = 2$ which is an integer and $S_{2m + 1} = S_1 = 1 \cdot \sqrt{17}$ and this holds true with the claim. Now for the inductive hypothesis, assume this claim holds for all $n < N$. We will show it also holds for $n = N$.

 Case 1: $N$ is even

 Then let $N = 2m$. Consider $S_{N - 1}$ and $S_{N - 2}$ so we get

 $S_{N - 1} = S_{2m - 1}$ and since $2m - 1 < N = 2m$, we can apply $S_{2m - 1} =  c_{m - 1} \cdot \sqrt{17}$ for some integer $c_{m - 1}$ to get $S_(N - 1) = c_{m - 1} \cdot \sqrt{17}$. Similarly,

 $S_{N - 2} = S_{2m - 2}$ and since $2m - 2 < N = 2m$, we can apply $S_{2m - 2}$ is an integer so

 $S_N = \sqrt{17} \cdot S_{N - 1} - 4 \cdot S_{N - 2} = \sqrt{17} \cdot c_{m - 1} \cdot \sqrt{17} - 4 \cdot \text{integer} = 17 \cdot c_{m - 1} - 4 \cdot \text{integer}$ which is obviously an integer. Thus, we've proven the claim for an even $N$.

 Case 2: $N$ is odd

 Then let $N = 2m + 1$. Consider $S_{N - 1}$ and $S_{N - 2}$ again to get

 $S_{N - 1} = S_{2m}$ and since $2m < N = 2m + 1$, we can apply $S_{2m}$ is an integer.

 $S_{N - 2} = S_{2m - 1}$ and again, as $2m - 1 < N = 2m + 1$, we can apply $S_{2m - 1} = c_{m - 1} \cdot \sqrt{17}$ for some integer $c_{m - 1}$ to get

 $S_N = \sqrt{17} \cdot S_{N - 1} - 4 \cdot S_{N - 2} = \sqrt{17} \cdot \text{integer} - 4 \cdot c_{m - 1} \cdot \sqrt{17} = \sqrt{17}(\text{integer} - 4 \cdot c_{m - 1}) = \sqrt{17} \cdot \text{another integer}$ which proves our claim for an odd $N$. $\square$

 So now we've proven that $k$ has to be even and since $n = 2k$, we have $n$ must be a multiple of $4$ which proves our first claim. $\square$

 So now we have $n = 4p$ and we need to find when $a^{\frac{n}{2}} + b^{\frac{n}{2}} = a^{2p} + b^{2p}$ is congruent to $9$ modulo $10$. Note that $a^2 = \frac{17 + 2\sqrt{17} + 1}{4} = \frac{9 + \sqrt{17}}{2}$ and similarly $b^2 = \frac{9 - \sqrt{17}}{2}$. Thus, we need to find when $(\frac{9 + \sqrt{17}}{2})^p + (\frac{9 - \sqrt{17}}{2})^p$ is an integer congruent to $9$ modulo $10$.

 We perform a similar substitution as before. Let $x = \frac{9 + \sqrt{17}}{2}, y = \frac{9 - \sqrt{17}}{2}$ so

 $x + y = 9, xy = \frac{81 - 17}{4} = 16$ so $x, y$ are the roots of the quadratic

 $t^2 - 9t + 16 = 0$

 Again, we introduce a characteristic polynomial, letting $x^k + y^k = P_k$ so that $P_k = 9 P_{k - 1} - 16 P_{k - 2}$. We manually compute $P_0 = x^0 + y^0 = 2, P_1 = x + y = 9$ so

 $P_2 = 9 \cdot P_1 - 16 \cdot P_0 = 9 \cdot 9 - 16 \cdot 2 = 81 - 32 = 49$. Also, modulo $10$ on all these terms so far gives $2, 9, 9$. We proceed.

 $P_3 = 9 \cdot P_2 - 16 \cdot P_1 = 9 \cdot 49 - 16 \cdot 9$ so modulo $10$ gives $-1 \cdot -1 - 16 \cdot -1 = 1 + 16 = 17$ which is $7$ modulo $10$. So we now have $2, 9, 9, 7$. So far, the only remainders of $9$ modulo $10$ were $k = 1, 2$ and $k = 0, 3$ do not work. Now, we suspect that all $k$ that are multiples of $3$ do not work. We need to prove this.

 Claim $P_{k}$ has a remainder of $9$ modulo $10$ whenever $k \equiv 1, 2 \mod 3$ but has a remainder of $7$ modulo $10$ whenever $k \equiv 0 \mod 3$

 Proof We will use induction. The base cases are $k = 1, 2, 3$ where we already established that these match the claim. Now, assume the claim holds for $1 \leq m < M$, for $M \geq 4$. We now show the claim holds for $m = M$.

 Note that we know $P_M = 9 \cdot P_{M - 1} - 16 \cdot P_{M - 2} \equiv -P_{M - 1} - 6P_{M - 2} \mod 10$.

 Case 1: $M \equiv 0 \mod 3$

 Then $M - 1 \equiv -1 \mod 3 \equiv 2 \mod 3$ and so $P_{M - 1} \equiv 9 \mod 10$ from our induction hypothesis since $M - 1 < M$. Similarly, $M - 2 \equiv -2 \mod 3 \equiv 1 \mod 3$ and so $P_{M - 2} \equiv 9 \mod 10$ from our induction hypothesis since $M - 2 < M$.

 So, $P_M \equiv -P_{M - 1} - 6P_{M - 2} \mod 10 \equiv -9 - 6(9) \mod 10 \equiv -9 - 54 \mod 10 \equiv -63 \mod 10 \equiv 7 \mod 10$ so the claim holds for $M \equiv 0 \mod 3$.

 Case 2: $M \equiv 1 \mod 3$

 Then $M - 1 \equiv 0 \mod 3$ so $P_{M - 1} \equiv 7 \mod 10$ from our induction hypothesis since $M - 1 < M$.Again $M - 2 \equiv -1 \mod 3 \equiv 2 \mod 3$ so $P_{M - 2} \equiv 9 \mod 10$ from our induction hypothesis since $M - 2 < M$. Thus,

 $P_M \equiv -P_{M - 1} - 6P_{M - 2} \mod 10 \equiv -7 - 6(9) \mod 10 \equiv -7 - 54 \mod 10 \equiv -61 \mod 10 \equiv 9 \mod 10$ so the claim holds for $M \equiv 1 \mod 3$.

 Case 3: $M \equiv 2 \mod 3$

 Then $M - 1 \equiv 1 \mod 3$ so $P_{M - 1} \equiv 9 \mod 10$ from our induction hypothesis since $M - 1 < M$.Also, $M - 2 \equiv 0 \mod 3$ so $P_{M - 1} \equiv 7 \mod 10$ from our induction hypothesis since $M - 2 < M$.

 Thus $P_M \equiv -P_{M - 1} - 6P_{M - 2} \mod 10 \equiv -9 - 6(7) \mod 10 \equiv -9 - 42 \mod 10 \equiv -51 \mod 10 \equiv 9 \mod 10$ so the claim holds for $M \equiv 2 \mod 3$.

 We've covered all the residues modulo $3$ and the claim holds for each one of the residues and thus we're done. $\square$

 We've proven all the main claims and now it is just computation. We have that every time $3 | k \implies 3 \cdot 4 | n \implies 12 | n$, we will get a remainder of $7$ modulo $10$. Every other $4 | n$ works. Thus, we count the number of multiples of $4$ less than or equal to $1000$ and subtract them from the number of multiples of $12$ less than or equal to $1000$ to get

 $\lfloor \frac{1000}{4} \rfloor - \lfloor \frac{1000}{12} \rfloor = 250 - 83 = \boxed{167}$



	\item (\textit{AIME II 2023}) A cube-shaped container has vertices $A,$ $B,$ $C,$ and $D,$ where $\overline{AB}$ and $\overline{CD}$ are parallel edges of the cube, and $\overline{AC}$ and $\overline{BD}$ are diagonals of faces of the cube, as shown. Vertex $A$ of the cube is set on a horizontal plane $\mathcal{P}$ so that the plane of the rectangle $ABDC$ is perpendicular to $\mathcal{P},$ vertex $B$ is $2$ meters above $\mathcal{P},$ vertex $C$ is $8$ meters above $\mathcal{P},$ and vertex $D$ is $10$ meters above $\mathcal{P}.$ The cube contains water whose surface is parallel to $\mathcal{P}$ at a height of $7$ meters above $\mathcal{P}.$ The volume of water is $\frac{m}{n}$ cubic meters, where $m$ and $n$ are relatively prime positive integers. Find $m+n.$

                    \begin{figure}[H]
                        \centering
                        \includegraphics[width=0.5\textwidth]{Problem/AIME_II_2023_14_Problem_Figure_1.png}
                    \end{figure}
                    

 \textbf{Answer: }751



\textbf{Solution 1}


                    \begin{figure}[H]
                        \centering
                        \includegraphics[width=0.5\textwidth]{Solution/AIME_II_2023_14_Solution_Figure_1.png}
                    \end{figure}
                    

 Let's first view the cube from a direction perpendicular to $ABDC$, as illustrated above. Let $x$ be the cube's side length. Since $\triangle CHA \sim \triangle AGB$, we have
 \[\frac{CA}{CH} = \frac{AB}{AG}.\]We know $AB = x$, $AG = \sqrt{x^2-2^2}$, $AC = \sqrt{2}x$, $CH = 8$. Plug them into the above equation, we get
 \[\frac{\sqrt{2}x}{8} = \frac{x}{\sqrt{x^2-2^2}}.\]Solving this we get the cube's side length $x = 6$, and $AC = 6\sqrt{2}.$

 Let $PQ$ be the water's surface, both $P$ and $Q$ are $7$ meters from $\mathcal P$. Notice that $C$ is $8$ meters from $\mathcal P$, this means
 \[CP = \frac{1}{8}CA = \frac{3\sqrt{2}}{4}.\]Similarly,
 \[DQ = \frac{3}{8}CA = \frac{9\sqrt{2}}{4}.\]

                    \begin{figure}[H]
                        \centering
                        \includegraphics[width=0.5\textwidth]{Solution/AIME_II_2023_14_Solution_Figure_2.png}
                    \end{figure}
                    

 Now, we realize that the 3D space inside the cube without water is a frustum, with $P$ on its smaller base and $Q$ on its larger base. To find its volume, all we need is to find the areas of both bases and the height, which is $x = 6$. To find the smaller base, let's move our viewpoint onto the plane $ABDC$ and view the cube from a direction parallel to $ABDC$, as shown above. The area of the smaller base is simply
 \[S_1 = CP^2 = \Bigl(\frac{3\sqrt{2}}{4}\Bigr)^2 = \frac{9}{8}.\]Similarly, the area of the larger base is
 \[S_2 = DQ^2 = \Bigl(\frac{9\sqrt{2}}{4}\Bigr)^2 = \frac{81}{8}.\]

 Finally, applying the formula for a frustum's volume,

 
 \[V = \frac{1}{3} \cdot x \cdot (S_1 + \sqrt{S_1S_2} + S_2) = \frac{1}{3} \cdot 6 \cdot \Bigl(\frac{9}{8} + \sqrt{\frac{9}{8}\cdot\frac{81}{8}} + \frac{81}{8}\Bigl) = \frac{117}{4}.\]

 The water's volume is thus
 \[6^3 - \frac{117}{4} = \frac{747}{4},\]giving $\boxed{751}$.

 



\textbf{Solution 2}


                    \begin{figure}[H]
                        \centering
                        \includegraphics[width=0.5\textwidth]{Solution/AIME_II_2023_14_Solution_Figure_3.png}
                    \end{figure}
                    Denote $h(X)$ the distance from point $X$ to $\mathcal{P},  h(A) = 0, h(B) = 2,$$h(C) = 8, h(D) = 10, h(G) = h(I) = h(H) = 7,  AB = a, AC = a \sqrt{2}.$

 Let slope $AB$ to $\mathcal{P}$ be $\alpha.$ Notation is shown in the diagram.
 \[\tan \alpha = \frac {\sin \alpha}{\cos \alpha} = \frac {h(B)}{AB}\cdot \frac {AC}{h(C)} = \frac {\sqrt{2}}{4} \implies a = 6.\]Let $S = GI \cap CD \implies h(S) = h(G) = 7.$
 \[h(C) – h(G) = 8 - 7 = 1, h(D)- h(I) = 10 - 7 = 3.\]
 \[h(E) = h(F) = \frac {h(D) +h(B)}{2} = 6 \implies\]
 \[\frac {DI}{DE} = \frac {h(D) - h(I)}{h(D)-h(E)} = \frac {3}{4} \implies DI = DH = \frac {9}{2}.\]

 Similarly $CG = \frac {3}{2} \implies SD = 9.$

 Let the volume without water be $V,$ volume of the pyramid $SCGJ$ be $U.$

 It is clear that $U + V = 27U = \frac {SD}{6} \cdot DI^2 = \frac {243}{8} \implies$$V = \frac {243 \cdot 26}{8 \cdot 27 } = \frac {117}{4} = 6^3 - \frac {747}{4}$  from which $\boxed{751}.$

 \textbf{vladimir.shelomovskii@gmail.com, vvsss}

 



\textbf{Solution 3}

We introduce a Cartesian coordinate system to the diagram.We put the origin at $A$. We let the $z$-components of $B$, $C$, $D$ be positive.We set the $x$-axis in a direction such that $B$ is on the $x-O-z$ plane.

 The coordinates of $A$, $B$, $C$ are $A = \left( 0, 0, 0 \right)$, $B = \left( x_B, 0 , 2 \right)$, $C = \left( x_C, y_C, 8 \right)$.

 Because $AB \perp AC$, $\overrightarrow{AB} \cdot \overrightarrow{AC} = 0$.Thus,
 \[ x_B x_C + 16 = 0 . \hspace{1cm} (1) \]

 Because $AC$ is a diagonal of a face, $AC^2 = 2 AB^2$.Thus,
 \[ x_C^2 + y_C^2 + 8^2 = 2 \left( x_B^2 + 2^2 \right) . \hspace{1cm} (2) \]

 Because plane $ABCD$ is perpendicular to plan $P$, $\hat z \cdot \left( \overrightarrow{AB} \times \overrightarrow{AC} \right) = 0$.Thus,
 \[ \begin{vmatrix} 0 & 0 & 1 \\ x_B & 0 & 2 \\ x_C & y_C & 8 \end{vmatrix} = 0 . \hspace{1cm} (3) \]

 Jointly solving (1), (2), (3), we get one solution $x_B = 4 \sqrt{2}$, $x_C = - 2 \sqrt{2}$, $y_C = 0$.Thus, the side length of the cube is $|AB| = \sqrt{x_B^2 + 2^2} = 6$.

 Denote by $P$ and $Q$ two vertices such that $AP$ and $AQ$ are two edges, and satisfy the right-hand rule that $\widehat{AB} \times \widehat{AP} = \widehat{AQ}$.Now, we compute the coordinates of $P$ and $Q$.

 Because $|AB| = 6$, we have $\overrightarrow{AP} \times \overrightarrow{AQ} = 6 \overrightarrow{AB}$, $\overrightarrow{AQ} \times \overrightarrow{AB} = 6 \overrightarrow{AP}$, $\overrightarrow{AB} \times \overrightarrow{AP} = 6 \overrightarrow{AQ}$.

 Hence,
 \begin{align*} \begin{bmatrix} \hat i & \hat j & \hat k \\ x_P & y_P & z_P \\ x_Q & y_Q & z_Q \end{bmatrix} & = 6 \left( 4 \sqrt{2} \hat i + 2 \hat k \right) , \\ \begin{vmatrix} \hat i & \hat j & \hat k \\ x_Q & y_Q & z_Q \\ 4 \sqrt{2} & 0 & 2 \end{vmatrix} & = 6 \left( x_P \hat i + y_P \hat j + z_P \hat k \right) , \\ \begin{vmatrix} \hat i & \hat j & \hat k \\ 4 \sqrt{2} & 0 & 2 \\ x_P & y_P & z_P \end{vmatrix} & = 6 \left( x_Q \hat i + y_Q \hat j + z_Q \hat k \right) . \end{align*}


 By solving these equations, we get$y_P^2 + y_Q^2 = 36 .$

 
In addition, we have $\overrightarrow{AC} = \overrightarrow{AP} + \overrightarrow{AQ}$.Thus, $P = \left( - \sqrt{2} , 3 \sqrt{2} , 4 \right)$, $Q = \left( - \sqrt{2} , - 3 \sqrt{2} , 4 \right)$.

 Therefore, the volume of the water is
 \begin{align*} V = & 6^3 \int_{u=0}^1 \int_{v=0}^1 \int_{w=0}^1 \mathbf 1 \left\{ z_B u + z_P v + z_Q w \leq 7 \right\} dw dv du \\ & = 6^3 \int_{u=0}^1 \int_{v=0}^1 \int_{w=0}^1 \mathbf 1 \left\{ 2 u + 4 v + 4 w \leq 7 \right\} dw dv du \\ & = 6^3 - 6^3 \int_{u=0}^1 \int_{v=0}^1 \int_{w=0}^1 \mathbf 1 \left\{ 2 u + 4 v + 4 w > 7 \right\} dw dv du . \end{align*}


 Define $u' = 1 - u$, $v' = 1 - v$, $w' = 1 - w$.Thus,
 \begin{align*} V & = 6^3 - 6^3 \int_{u=0}^1 \int_{v=0}^1 \int_{w=0}^1 \mathbf 1 \left\{ 2 u' + 4 v' + 4 w' < 3 \right\} dw dv du \\ & = 6^3 - 6^3 \int_{u'=0}^1 \left( \int_{v'=0}^1 \int_{w'=0}^1  \mathbf 1 \left\{ v' + w' < \frac{3}{4} - \frac{u'}{2} \right\} dw' dv' \right)  du' \\ & = 6^3 - 6^3 \int_{u'=0}^1  \frac{1}{2} \left( \frac{3}{4} - \frac{u'}{2} \right)^2 du' . \end{align*}


 Define $u'' = \frac{3}{4} - \frac{u'}{2}$.Thus,
 \begin{align*} V & = 6^3 - 6^3 \int_{u'' = 1/4}^{3/4} \left( u'' \right)^2 du'' \\ & = 6^3 - 6^3 \frac{1}{3} \left( \left(\frac{3}{4}\right)^3 - \left(\frac{1}{4}\right)^3 \right) \\ & = 216 - \frac{117}{4} \\ & = \frac{747}{4} . \end{align*}


 Therefore, the answer is $747 + 4 = \boxed{\textbf{(751) }}$.

 


	\item (\textit{AIME II 2023}) For each positive integer $n$ let $a_n$ be the least positive integer multiple of $23$ such that $a_n \equiv 1 \pmod{2^n}.$ Find the number of positive integers $n$ less than or equal to $1000$ that satisfy $a_n = a_{n+1}.$

 \textbf{Answer: }363



\textbf{Solution 1}

Denote $a_n = 23 b_n$.Thus, for each $n$, we need to find smallest positive integer $k_n$, such that
 \[ 23 b_n = 2^n k_n + 1 . \]

 Thus, we need to find smallest $k_n$, such that
 \[ 2^n k_n \equiv - 1 \pmod{23} . \]

 Now, we find the smallest $m$, such that $2^m \equiv 1 \pmod{23}$.By Fermat's Theorem, we must have $m | \phi \left( 23 \right)$. That is, $m | 22$.We find $m = 11$.

 Therefore, for each $n$, we need to find smallest $k_n$, such that
 \[ 2^{{\textrm Rem} \left( n , 11 \right)} k_n \equiv - 1 \pmod{23} . \]

 We have the following results:

 If \({\textrm Rem} \left( n , 11 \right) = 0\), then \(k_n = 22\) and \(b_n = 1\).

 If \({\textrm Rem} \left( n , 11 \right) = 1\), then \(k_n = 11\) and \(b_n = 1\).

 If \({\textrm Rem} \left( n , 11 \right) = 2\), then \(k_n = 17\) and \(b_n = 3\).

 If \({\textrm Rem} \left( n , 11 \right) = 3\), then \(k_n = 20\) and \(b_n = 7\).

 If \({\textrm Rem} \left( n , 11 \right) = 4\), then \(k_n = 10\) and \(b_n = 7\).

 If \({\textrm Rem} \left( n , 11 \right) = 5\), then \(k_n = 5\) and \(b_n = 7\).

 If \({\textrm Rem} \left( n , 11 \right) = 6\), then \(k_n = 14\) and \(b_n = 39\).

 If \({\textrm Rem} \left( n , 11 \right) = 7\), then \(k_n = 7\) and \(b_n = 39\).

 If \({\textrm Rem} \left( n , 11 \right) = 8\), then \(k_n = 15\) and \(b_n = 167\).

 If \({\textrm Rem} \left( n , 11 \right) = 9\), then \(k_n = 19\) and \(b_n = 423\).

 If \({\textrm Rem} \left( n , 11 \right) = 10\), then \(k_n = 21\) and \(b_n = 935\).

 Therefore, in each cycle, $n \in \left\{ 11 l , 11l + 1 , \cdots , 11l + 10 \right\}$, we have $n = 11l$, $11l + 3$, $11l + 4$, $11l + 6$, such that $b_n = b_{n+1}$. That is, $a_n = a_{n+1}$.At the boundary of two consecutive cycles, $b_{11L + 10} \neq b_{11 \left(l + 1 \right)}$.

 We have $1000 = 90 \cdot 11 + 10$.Therefore, the number of feasible $n$ is $91 \cdot 4 - 1 = \boxed{\textbf{(363) }}$.

 



\textbf{Solution 2}

Observe that if $a_{n-1} - 1$ is divisible by $2^n$, $a_n = a_{n-1}$. If not, $a_n = a_{n-1} + 23 \cdot 2^{n-1}$.

 This encourages us to let $b_n = \frac{a_n - 1}{2^n}$. Rewriting the above equations, we have
 \[b_n = \begin{cases} \frac{b_{n-1}}{2} & \text{if } 2 \text{ } \vert \text{ } b_{n-1} \\ \frac{b_{n-1}+23}{2} &\text{if } 2 \not\vert \text{ } b_{n-1} \end{cases}\]The first few values of $b_n$ are $11, 17, 20, 10, 5, 14, 7, 15, 19, 21,$ and $22$. We notice that $b_{12} = b_1 = 11$, and thus the sequence is periodic with period $11$.

 Note that $a_n = a_{n+1}$ if and only if $b_n$ is even. This occurs when $n$ is congruent to $0, 3, 4$ or $6$ mod $11$, giving four solutions for each period.

 From $1$ to $1001$ (which is $91 \times 11$), there are $91 \times 4 = 364$ values of $n$. We subtract $1$ from the total since $1001$ satisfies the criteria but is greater than $1000$ to get a final answer of $\fbox{363}$ .~\href{https://artofproblemsolving.com/wiki/index.php?title=User:Bxiao31415}{Bxiao31415}

 (small changes by bobjoebilly and \href{https://artofproblemsolving.com/wiki/index.php?title=User:Iraevid13}{IraeVid13})

 



\textbf{Solution 3 (Similar to solution 2 but more explanation)}

Let $a_n = 23b_n$. Note that if
 \[23 b_{n+1} \equiv 1 \pmod{2^{n+1}}\]Then
 \[23 b_{n+1} \equiv 1 \pmod{2^{n}}\]Also
 \[23 b_n \equiv 1 \pmod{2^n}\]Therefore
 \[b_n \equiv b_{n+1} \equiv 23^{-1} \pmod{2^n}\]Then
 \[b_{n+1} \equiv b_n, b_n + 2^n \pmod{2^{n+1}}\]So
 \[b_{n+1} = b_n, b_n + 2^n\]Since $0 \le b_n < 2^n$ and $0 \le b_{n+1} < 2^{n+1}$ as $a_n$ is the \textit{least} positive integer multiple of 23.

 Now suppose $b_{n+1} = b_n$. Define $q_n$ to be the quotient of $23b_n$ divided by $2^n$. Then
 \[23b_n = 2^n q_n + 1 \text{ and } 23b_{n+1} = 23b_n = 2^{n+1} q_{n+1} + 1 = 2^n q_n + 1 \implies q_{n+1} = \frac{q_n}{2}\]Furthermore if quotient $q_n$ is even then
 \[23b_n = 2^n q_n +1 = 2^{n+1} \frac{q_n}{2} +1\]Therefore $b_{n+1} = b_n$ if and only if $q_n$ is even. And, if this is true, then $q_{n+1} = \frac{q_n}{2}$. Next, if $q_n$ is odd, we must have $b_{n+1} = b_n + 2^n$. Solving for $q_{n+1}$, we have
 \[23b_{n+1} = 2^{n+1} q_{n+1} + 1 \implies 23b_n + 23 \cdot 2^n = 2^{n+1} q_{n+1} + 1 \implies 2^n q_n + 1 + 23 = 2^{n+1} q_{n+1} + 1 \implies q_{n+1} = \frac{q_n + 1}{2} + 11\]Therefore, if $q_n$ is odd, $q_{n+1} = \frac{q_n + 1}{2} + 11$. In sum, our recursion is 
 \[q_n = \begin{cases} \frac{q_{n-1}}{2} & \text{if } 2 \text{ } \vert \text{ } q_{n-1} \\ \frac{q_{n-1}+1}{2} + 11 &\text{if } 2 \not\vert \text{ } q_{n-1} \end{cases}\]Finally, let us list out $q_n$ to find a pattern. Because $a_1 = 23$, $q_1 = 11$. Through our recursion, we continue like so:
 \[q_1 = 11, q_2 = 17, q_2 = 20, q_3 = 10, q_4 = 5, q_6 = 14, q_7 = 7, q_8 = 15, q_9 = 19, q_{10} = 21, q_{11} = 22, q_{12} = 11, \dots\]Therefore $q_n$ repeats with cycle length $11$. Since $a_{n+1} = a_n$ if and only if $q_n$ is even, in each cycle, we have 4 satisfactory values of $n$. There are $\frac{1000 - 10}{11} = 90$ complete cycles. There are 3 extra values in the last incomplete cycle. Therefore we obtain $90 \cdot 4 + 3 = \fbox{363}$.

 



\textbf{Solution 4 (Binary Interpretation, Computer Scientists' Playground)}

We first check that $\gcd(23, 2^n) = 1$ hence we are always seeking a unique modular inverse of $23$, $b_n$, such that $a_n \equiv 23b_n \equiv 1 \mod{2^n}$.

 
Now that we know that $b_n$ is unique, we proceed to recast this problem in binary. This is convenient because $x \mod{2^n}$ is simply the last $n$-bits of $x$ in binary, and if $x \equiv 1 \mod{2^n}$, it means that of the last $n$ bits of $x$, only the rightmost bit (henceforth $0$th bit) is $1$.

 Also, multiplication in binary can be thought of as adding shifted copies of the multiplicand. For example:

 
 \begin{align} 10111_2 \times 1011_2 &= 10111_2 \times (1000_2 + 10_2 + 1_2) \\               &= 10111000_2 + 101110_2 + 10111_2 \\               &= 11111101_2 \end{align}

 Now note $23 = 10111_2$, and recall that our objective is to progressively zero out the $n$ leftmost bits of $a_n = 10111_2 \times b_n$ except for the $0$th bit. 

 Write $b_n = \underline{c_{n-1}\cdots c_2c_1c_0}_2$, we note that $c_0$ uniquely defines the $0$th bit of $a_n$, and once we determine $c_0$, $c_1$ uniquely determines the $1$st bit of $a_n$, so on and so forth. 

 For example, $c_0 = 1$ satisfies $a_1 \equiv10111_2 \times 1_2 \equiv 1 \mod{10_2}$Next, we note that the second bit of $a_1$ is $1$, so we must also have $c_1 = 1$ in order to zero it out, giving

 
 \[a_2 \equiv 10111_2 \times 11_2 \equiv 101110_2 + a_1 \equiv 1000101_2 \equiv 01_2 \mod{100_2}\]
$a_{n+1} = a_{n}$ happens precisely when $c_n = 0$. In fact we can see this in action by working out $a_3$. Note that $a_2$ has 1 on the $2$nd bit, so we must choose $c_2 = 1$. This gives 

 
 \[a_3 \equiv 10111_2 \times 111_2 \equiv 1011100_2 + a_2 \equiv 10100001_2 \equiv 001_2 \mod{1000_2}\]

 Note that since the $3$rd and $4$th bit are $0$, $c_3 = c_4 = 0$, and this gives $a_3 = a_4 = a_5$.

 
It may seem that this process will take forever, but note that $23 = 10111_2$ has $4$ bits behind the leading digit, and in the worst case, the leading digits of $a_n$ will have a cycle length of at most $16$. In fact, we find that the cycle length is $11$, and in the process found that $a_3 = a_4 = a_5$, $a_6 = a_7$, and $a_{11} = a_{12}$.

 Since we have $90$ complete cycles of length $11$, and the last partial cycle yields $a_{993} = a_{994} = a_{995}$ and $a_{996} = a_{997}$, we have a total of $90 \times 4 + 3 = \boxed{363}$ values of $n \le 1000$ such that $a_n = a_{n+1}$




\textbf{Solution 5 (kind of sketchy, intuitive)}

Let $a_n=23k$. Note that each time $n$ increases, either $a_n$ remains the same or we add $2^n$ to $k$, which motivates binary. The condition is equivalent to $23k \equiv 1 \pmod{2^n} \implies k \equiv \frac{1}{23} \pmod{2^n}$. Then, consider the binary expansion of $\frac{1}{23}$; it can be shown that this is equal to $0.\overline{00001011001}$. 

 


 (The way you get this is basically setting $\frac{x}{2^k} + \frac{x}{2^{2k}} \ldots = \frac{1}{23} \implies \frac{x}{2^k-1} = \frac{1}{23}$. Then, $23 \mid 2^k-1$, and the minimal such $k$ which satisfies this is $k=11$. This can be obtained by noting that $k \mid \varphi(23) = 22$. Then, cross multiply to solve for $x$.) 

 


 Then, if $a_n = a_{n+1}$, the $n$th digit in the binary expansion is a $1$ because of reasons. Of the first 1000 digits in the binary expansion this happens $90 \cdot 4 + 3 = \boxed{363}$ times.

 


\end{enumerate}